\documentclass[11pt,oneside]{article}

\usepackage[utf8]{inputenc}
\usepackage[T1]{fontenc}
\usepackage{lmodern}
\usepackage{microtype}
\usepackage{geometry}
\geometry{margin=1in,headheight=14pt}
\usepackage{amsmath,amssymb,amsthm,mathtools}
\usepackage{booktabs,longtable,tabularx,array,multirow}
\usepackage{enumitem}
\usepackage{xcolor}
\usepackage{xurl}
\usepackage{hyperref}
\usepackage[nameinlink,noabbrev]{cleveref}
\usepackage{listings}
\usepackage{fancyhdr}
\usepackage{titlesec}
\usepackage{needspace}

\definecolor{leanblue}{RGB}{24,91,145}
\definecolor{newgreen}{RGB}{15,105,68}
\definecolor{computeviolet}{RGB}{98,55,145}
\definecolor{classicalgray}{RGB}{85,92,99}
\definecolor{codegray}{RGB}{248,248,248}
\definecolor{rulegray}{RGB}{92,101,111}
\definecolor{warmbox}{RGB}{249,247,239}

\hypersetup{
  colorlinks=true,
  linkcolor=leanblue,
  citecolor=newgreen,
  urlcolor=leanblue,
  pdftitle={The Alaoglu--Erdos Integer-Exponent Conjecture: Further Determinant Methods},
  pdfauthor={Research report for the ProveIt project},
  pdfsubject={Arbitrary-node repulsion, semigroup determinants, exact conditional structure, and research directions}
}
\urlstyle{same}

\pagestyle{fancy}
\fancyhf{}
\lhead{Alaoglu--Erd\H{o}s integer-exponent conjecture}
\rhead{Further research report}
\cfoot{\thepage}
\renewcommand{\headrulewidth}{0.4pt}

\titleformat{\section}{\Large\bfseries}{\thesection}{0.75em}{}
\titleformat{\subsection}{\large\bfseries}{\thesubsection}{0.75em}{}
\titleformat{\subsubsection}{\normalsize\bfseries}{\thesubsubsection}{0.75em}{}

\setlist[itemize]{leftmargin=2em,itemsep=0.28em,topsep=0.4em}
\setlist[enumerate]{leftmargin=2.2em,itemsep=0.28em,topsep=0.4em}
\setlist[description]{leftmargin=2.9cm,labelsep=0.65cm,itemsep=0.35em,topsep=0.4em}
\setlength{\parindent}{0pt}
\setlength{\parskip}{0.55em}
\setlength{\emergencystretch}{3.5em}

\newtheorem{theorem}{Theorem}[section]
\newtheorem{proposition}[theorem]{Proposition}
\newtheorem{lemma}[theorem]{Lemma}
\newtheorem{corollary}[theorem]{Corollary}
\newtheorem{conjecture}[theorem]{Conjecture}
\newtheorem{question}[theorem]{Question}
\theoremstyle{definition}
\newtheorem{definition}[theorem]{Definition}
\newtheorem{observation}[theorem]{Observation}
\theoremstyle{remark}
\newtheorem{remark}[theorem]{Remark}

\newcommand{\N}{\mathbb N}
\newcommand{\Nzero}{\mathbb N_0}
\newcommand{\Z}{\mathbb Z}
\newcommand{\Q}{\mathbb Q}
\newcommand{\R}{\mathbb R}
\newcommand{\C}{\mathbb C}
\newcommand{\Abar}{\overline{\mathbb Q}}
\newcommand{\thetaq}{\vartheta}
\newcommand{\Sol}{\mathcal S}
\newcommand{\Cand}{\mathcal C}
\newcommand{\LambdaTheta}{\Lambda_{\thetaq}}
\newcommand{\Zthree}{\Z[1/3]}
\newcommand{\Ztwo}{\Z[1/2]}
\newcommand{\fract}[1]{\left\{#1\right\}}
\newcommand{\distZ}[1]{\left\lVert #1\right\rVert_{\Z}}
\newcommand{\vpp}[2]{v_{#1}\!\left(#2\right)}
\newcommand{\fall}[2]{\left(#1\right)_{\!\underline{#2}}}
\newcommand{\lean}[1]{\texttt{\detokenize{#1}}}
\newcommand{\commit}[1]{\texttt{#1}}
\newcommand{\statuslean}{\textcolor{leanblue}{\textbf{[repository kernel-verified Lean]}}}
\newcommand{\statusnew}{\textcolor{newgreen}{\textbf{[paper proof in this report]}}}
\newcommand{\statuscomp}{\textcolor{computeviolet}{\textbf{[interval computation]}}}
\newcommand{\statusknown}{\textcolor{classicalgray}{\textbf{[classical/open background]}}}

\newenvironment{statusbox}[1]
{\begin{center}\begin{minipage}{0.94\textwidth}\raggedright\hbadness=10000\hrule\vspace{0.45em}\textbf{Status:} #1\par\vspace{0.35em}}
{\vspace{0.35em}\hrule\end{minipage}\end{center}}

\newenvironment{resultbox}
{\begin{center}\begin{minipage}{0.94\textwidth}\hrule\vspace{0.45em}}
{\vspace{0.35em}\hrule\end{minipage}\end{center}}

\lstdefinestyle{leanstyle}{
  basicstyle=\ttfamily\footnotesize,
  backgroundcolor=\color{codegray},
  frame=single,
  rulecolor=\color{rulegray},
  showstringspaces=false,
  breaklines=true,
  columns=fullflexible,
  keepspaces=true,
  tabsize=2,
  captionpos=b
}

\begin{document}
\pagenumbering{gobble}

\begin{titlepage}
\centering
\vspace*{0.75cm}
{\Huge\bfseries The Alaoglu--Erd\H{o}s\par}
\vspace{0.2cm}
{\Huge\bfseries Integer-Exponent Conjecture\par}
\vspace{0.65cm}
{\LARGE Further Determinant Methods,\par}
\vspace{0.12cm}
{\LARGE Arbitrary-Node Repulsion,\par}
\vspace{0.12cm}
{\LARGE and the Remaining Four-Exponentials Barrier\par}
\vspace{0.95cm}
{\large A repository-grounded research report based on\par}
\vspace{0.12cm}
{\large \url{https://github.com/VladimirReshetnikov/ProveIt}\par}
\vspace{0.35cm}
{\normalsize pinned at commit\par}
{\normalsize \commit{5d3982ceb12644fdd3499569d9db5f62f9a31dc6}\par}
\vfill
\begin{minipage}{0.9\textwidth}
\small
\textbf{Bottom line.}
No unconditional proof of the conjecture is obtained.  The current ProveIt corpus already
pushes the conditional structure remarkably far: if a counterexample exists, then all
solutions form the free additive monoid $\Nzero+\Nzero\beta$ generated by $1$ and one least
nonintegral exponent $\beta$, and all natural candidate bases form the free multiplicative
monoid generated by $2$ and $M=2^\beta$.

This report adds two paper-level determinant methods.  First, arbitrary candidate tuples on
the lattice curve $y=x^{\thetaq}$ satisfy an explicit Vandermonde repulsion inequality; this
strictly extends the repository's arithmetic-progression obstructions to completely irregular
nodes.  Second, the integrality of every mixed monomial
$m^{a+b\thetaq}=m^a(m^{\thetaq})^b$ gives a hierarchy of positive integer generalized
Vandermonde determinants.  It yields an independent, explicit
$O((\log X)^2)$ bound for candidates in $[X,2X]$, without using the prime-valuation rank theorem.
The method also diagnoses its own failure: a hypothetical counterexample supplies only
$\asymp\log X$ points in such a block, while the archimedean determinant needs
$\asymp(\log X)^2$ points to become nearly scale-forcing.  Closing precisely this one-logarithm
deficit appears to require a genuinely new ingredient, most plausibly $p$-adic divisibility,
a denominator-sensitive counting theorem, or a new algebraic exponential.
\end{minipage}
\vfill
{\large August 20, 2026\par}
\end{titlepage}

\pagenumbering{roman}

\begin{abstract}
Let $\thetaq=\log 3/\log 2$.  The Alaoglu--Erd\H{o}s conjecture is equivalent to the assertion
that the only positive integers $m$ for which $m^{\thetaq}$ is an integer are the powers of
$2$.  The pinned ProveIt development contains kernel-checked zero-free regions,
transcendence restrictions, multiplicative-rank and density bounds, a canonical primitive odd
core, a localized-radical reformulation, exact radical degree, arbitrary-order
arithmetic-progression obstructions, and --- conditional on failure --- a complete
classification
\[
   \Sol=\{n+k\beta:n,k\in\Nzero\}
\]
with unique coordinates.  It also gives exact dyadic-block counts and a sharp quadratic
cumulative asymptotic.  These results reduce the conjecture to a very narrow special case of
the four exponentials problem: no primitive odd $w>1$ should have a positive power of
$w^{\thetaq}$ in $\Z[1/3]$.

We attempt several routes to the remaining obstruction and make their failure points explicit.
The principal new theorem is an arbitrary-node interpolation determinant.  If
$1\le m_0<\cdots<m_r$ and every $m_i^{\thetaq}$ is integral, then for $r\ge2$
\[
 \prod_{0\le i<j\le r}(m_j-m_i)
 \ge
 \frac{r!}{|\fall{\thetaq}{r}|}\,m_0^{r-\thetaq}.
\]
For three candidates in $[N,N+H]$ this gives
\[
 H\ge
 \left(\frac{8}{\thetaq(\thetaq-1)}\right)^{1/3}
 N^{(2-\thetaq)/3}
 =2.05107\ldots\,N^{0.138345\ldots}.
\]
The order-$4$ member gives the best exponent among the ordinary-power family: every interval
shorter than $1.46022\,N^{0.241503\ldots}$ contains at most four candidates.

A stronger hierarchy uses the ordered semigroup
$\LambdaTheta=\{a+b\thetaq:a,b\in\Nzero\}$.  For its first $r+1$ exponents
$\lambda_0<\cdots<\lambda_r$, every candidate tuple produces the strictly positive integer
\[
  \det\bigl(m_i^{\lambda_j}\bigr)_{0\le i,j\le r}.
\]
A divided-difference estimate gives, with
$R=r(r+1)/2$ and $\Sigma_r=\sum_{j=0}^r\lambda_j$,
\[
 H\ge ((r+1)!)^{-1/R}2^{-1-2\Sigma_r/R}
       N^{1-\Sigma_r/R}
\]
whenever $r+1$ candidates lie in $[N,N+H]\subseteq[N,2N]$.
Since
$\Sigma_r/R=(4/3)\sqrt{2\thetaq/r}+O(r^{-1})$, the exponent tends to $1$.
An entirely explicit corollary is
\[
 \#\bigl(\Cand\cap[X,2X]\bigr)
 \le 10000\max\{1,\log X\}^{2}.
\]
This is independent of the repository's valuation-rank proof and is suitable for Lean
formalization.  It still does not contradict the exact conditional monoid, whose block count
is only $\asymp\log X$.  The report ends with new equivalent formulations, a prioritized
research program, and a module-by-module formalization blueprint.
\end{abstract}

\tableofcontents
\clearpage
\pagenumbering{arabic}

\section{Executive verdict}
\label{sec:executive}

\subsection{What was attempted}

The investigation began from the current combined report and the Lean modules imported by
\lean{ExponentialIdentities.lean} at commit
\commit{5d3982ceb12644fdd3499569d9db5f62f9a31dc6}.  The main lines pursued were:

\begin{enumerate}
\item an exact reduction to the integer curve $m\mapsto m^{\thetaq}$ and to the logarithmic
$2\times2$ determinant behind four exponentials;
\item attempts to combine the canonical primitive-generator theorem with algebraic-number
and radical-field arguments;
\item rational approximation to a hypothetical least generator, including simultaneous
integer lower bounds at bases $2$ and $3$;
\item density and shrinking-target arguments for the fractional parts of its multiples;
\item finite differences beyond arithmetic progressions;
\item interpolation determinants on arbitrary candidate nodes;
\item generalized Vandermonde determinants using the full semigroup of integral mixed
monomials $m^{a+b\thetaq}$;
\item possible $p$-adic, $S$-integral, and auxiliary-base amplifications.
\end{enumerate}

The first five routes reach known barriers.  The sixth and seventh produce new unconditional
partial results.  None closes the final transcendence-theoretic gap.

\subsection{The strongest new conclusions}

\begin{resultbox}
\textbf{New theorem A: arbitrary-node repulsion.}
For every $r\ge2$, any $r+1$ distinct natural candidates obey a quantitative product-of-gaps
lower bound.  No arithmetic-progression hypothesis is needed.  This gives explicit local
packing bounds and a best ordinary-power exponent $0.241503\ldots$ at determinant order $4$.
\end{resultbox}

\begin{resultbox}
\textbf{New theorem B: semigroup determinant hierarchy.}
The mixed exponent semigroup $\Nzero+\Nzero\thetaq$ supplies arbitrarily large positive
integer generalized Vandermonde determinants.  Their archimedean upper bounds yield a
near-linear repulsion exponent $1-O(r^{-1/2})$ and, unconditionally,
\[
 \#(\Cand\cap[X,2X])\ll(\log X)^2.
\]
An explicit constant $10000$ is proved below.  This route is independent of the repository's
prime-valuation rank theorem.
\end{resultbox}

\begin{resultbox}
\textbf{New diagnosis: the one-logarithm deficit.}
To make the semigroup determinant nearly force a gap comparable with $X$, one needs
$r\asymp(\log X)^2$ candidate nodes.  Conditional on a counterexample, the exact ProveIt
classification supplies only $\asymp\log X$ candidates in $[X,2X]$.  Thus purely
archimedean interpolation, at its present strength, misses the conditional population by one
power of $\log X$.  A successful refinement must either reduce the determinant's node
requirement to $O(\!\log X)$ or extract additional nonarchimedean size from the same nodes.
\end{resultbox}

\subsection{Why this does not amount to a proof}

The conjecture is already a special integral instance of the real four exponentials
conjecture.  In the natural $2\times2$ matrix
\[
 \begin{pmatrix}
  1\cdot\log2 & 1\cdot\log3\\
  x\log2 & x\log3
 \end{pmatrix},
\]
the four exponentials are $2,3,2^x,3^x$.  For a nonintegral solution, both row and column
pairs are rationally linearly independent, yet all four exponentials are algebraic.  The
conclusion needed here is therefore exactly what four exponentials predicts and what the
known six exponentials theorem does not supply in a $2\times2$ configuration.

The additional integrality hypotheses are stronger than mere algebraicity and do yield the
new determinant inequalities.  They do not, by themselves, produce enough nodes or enough
lower-bound divisibility to cross the four-exponentials threshold.  The full conjecture should
therefore remain marked open in this report.

\section{Evidence boundary and repository provenance}
\label{sec:provenance}

\subsection{Pinned state}

The report is based on the repository state at
\[
 \commit{5d3982ceb12644fdd3499569d9db5f62f9a31dc6},
\]
whose merge message records the reconciliation of the parallel conditional-structure and
asymptotic formalizations.  The source file specifically requested by the user is
\begin{center}
\lean{Analysis/ExponentialIdentities/Article/alaoglu-erdos-combined-report/}
\lean{alaoglu-erdos-combined-report.tex}.
\end{center}
The authoritative formal import surface is \lean{ExponentialIdentities.lean}.

This report inspected the cited Lean sources and theorem statements.  It did not rerun the
entire Lean build locally.  Consequently, \statuslean{} means that the proof is present on the
pinned repository import surface and is reported there as accepted by the Lean kernel; it is
not a claim that a separate fresh build was performed while preparing this document.

\subsection{Status labels}

\begin{description}
\item[\statuslean]
A theorem is implemented in the pinned Lean corpus and appears on the aggregate import
surface.  The exact identifier is supplied whenever it is central.

\item[\statusnew]
A conventional mathematical proof is given in this report.  It is written with formalization
in mind but has not yet been translated into Lean.

\item[\statuscomp]
A finite interval computation described in the combined report.  It has a stated software
trust boundary and is not a kernel certificate.

\item[\statusknown]
A standard theorem, an established open conjecture, or bibliographic background.
\end{description}

\subsection{Current formal frontier}

\begin{longtable}{>{\raggedright\arraybackslash}p{0.32\textwidth}
                  >{\raggedright\arraybackslash}p{0.22\textwidth}
                  >{\raggedright\arraybackslash}p{0.38\textwidth}}
\toprule
Result & Status & Principal Lean identifier or source \\
\midrule
\endhead
Three-base theorem $2^x,3^x,5^x\in\Z\Rightarrow x\in\Z$
& \statuslean & \lean{IntegerExponent.lean}; interpolation determinant \\
Every nonintegral solution is transcendental
& \statuslean & \path{transcendental_of_not_integer_of_two_three_rpow_integer} \\
$\thetaq=\log3/\log2$ is transcendental
& \statuslean & \path{transcendental_logThreeDivLogTwo} \\
No nonintegral solution with $2^x<4096$
& \statuslean & \path{FiniteCheck4096.lean} \\
Independent exclusion with $2^x<2^{20}$
& \statuscomp & Combined report interval scan; not a Lean proof \\
Four-exponentials implication
& \statuslean & \path{alaogluErdosConjecture_of_realFourExponentialsConjecture} \\
Canonical primitive odd core and exact least generator
& \statuslean & \path{exists_canonical_primitiveGenerator_of_not_alaogluErdosConjecture} \\
Exact solution monoid $\Nzero+\Nzero\beta$
& \statuslean & Same primitive-generator theorem \\
Exact candidate monoid $\{2^nM^k\}$
& \statuslean & Same primitive-generator theorem \\
Exact count in $[2^T,2^{T+1})$
& \statuslean & \path{twoBaseNaturalCandidatesInDyadicBlock_card} \\
Sharp quadratic cumulative asymptotic
& \statuslean & \path{tendsto_twoBaseNaturalCandidatesBelow_two_pow_div_sq} \\
Growth equivalence for the conjecture
& \statuslean & \path{alaogluErdosConjecture_iff_candidate_count_below_two_pow_subquadratic} \\
Vanishing block Fourier modes
& \statuslean & \path{exists_primitiveGenerator_with_vanishing_block_fourier_modes} \\
Localized radical equivalence
& \statuslean & \path{alaogluErdosConjecture_iff_odd_rpow_not_localized} \\
Primitive localized radical equivalence
& \statuslean & \path{alaogluErdosConjecture_iff_primitiveOddLocalizedRadicalExclusion} \\
Exact binomial minimal polynomial and radical degree
& \statuslean & \path{oddCoreRpow_exact_radical_degree} \\
Simultaneous primitive output normalization
& \statuslean & \path{exists_simultaneous_primitive_output_normalization} \\
Arbitrary-order arithmetic-progression obstruction
& \statuslean & \path{not_arithmetic_progression_twoBaseNaturalCandidates_of_curvature} \\
Arbitrary-node Vandermonde repulsion
& \statusnew & \Cref{thm:arbitrary-node} \\
Semigroup generalized-Vandermonde hierarchy
& \statusnew & \Cref{thm:semigroup-integer-det,thm:semigroup-gap} \\
Explicit dyadic $O((\log X)^2)$ bound by that hierarchy
& \statusnew & \Cref{cor:semigroup-log-square} \\
Full Alaoglu--Erd\H{o}s conjecture
& Open & Equivalent radical and four-exponentials barriers below \\
\bottomrule
\end{longtable}


\subsection{Historical placement}

Alaoglu and Erd\H{o}s isolated the two-prime assertion in their 1944 study of highly
composite and related numbers.  They regarded it as very likely but could not prove it, and
reported a communicated three-prime result of Siegel.  In modern language, the three-base
statement is supplied by the six exponentials theorem, whereas the two-base statement sits at
the unresolved $2\times2$ threshold of the four exponentials conjecture
\cite{AE1944,Waldschmidt2000,Waldschmidt2005}.  Thus the elementary-looking integrality
question is historically one of the canonical concrete tests of the gap between six and four
exponentials, not an isolated recreational puzzle.

\section{The problem in its most useful forms}
\label{sec:setup}

\subsection{Exponent form and candidate form}

\begin{conjecture}[Alaoglu--Erd\H{o}s]
\label{conj:AE}
If $x\in\R$ and $2^x,3^x\in\Z$, then $x\in\Z$.
\end{conjecture}

Because $2^x$ and $3^x$ are positive, ``integer'' can be replaced by ``positive natural
number.''  A solution cannot have $x<0$, since then $0<2^x<1$.  Put
\[
  \thetaq:=\frac{\log3}{\log2}=1.584962500721156\ldots
\]
and define the natural candidate set
\[
  \Cand:=\{m\in\N:m>0,\ m^{\thetaq}\in\N\}.
\]
If $m=2^x$, then
\[
  m^{\thetaq}=(2^x)^{\log3/\log2}=3^x.
\]
Conversely, for $m>0$, the exponent $x=\log_2m$ satisfies $2^x=m$ and
$3^x=m^{\thetaq}$.  Therefore
\begin{equation}
 \boxed{\quad\text{Conjecture~\ref{conj:AE}}
 \iff \Cand=\{2^n:n\in\Nzero\}.\quad}
 \label{eq:candidate-equivalence}
\end{equation}
This is the form used throughout the new determinant arguments.

\subsection{The logarithmic determinant}

Writing
\[
  M=2^x,\qquad A=3^x,
\]
a solution is equivalently a positive integer pair satisfying
\begin{equation}
  \log M\,\log3-\log A\,\log2=0.
  \label{eq:log-det}
\end{equation}
The conjecture says that all positive integer solutions of \eqref{eq:log-det} are
\[
  (M,A)=(2^n,3^n),\qquad n\in\Nzero.
\]
Thus the problem is a vanishing $2\times2$ determinant of logarithms of algebraic numbers.
This is not merely an analogy with four exponentials; it is its exact local geometry.

\subsection{One-parameter subgroup form}

Let
\[
  \Gamma=\{(2^t,3^t):t\in\R\}\subset\R_{>0}^2.
\]
Then the conjecture is
\[
  \Gamma\cap\N^2=\{(2^n,3^n):n\in\Nzero\}.
\]
The curve is a real one-parameter subgroup of the multiplicative torus
$\mathbb G_m^2(\R)$.  The difficulty is that its direction vector
$(\log2,\log3)$ is transcendental rather than algebraic, so standard algebraic-subgroup
intersection theorems do not apply directly.

\subsection{Basic unconditional restrictions}

The pinned corpus proves the following.

\begin{itemize}
\item A nonintegral $x$ with $2^x\in\Z$ is irrational, and in fact transcendental by
Gelfond--Schneider.  Hence every counterexample is transcendental.  \statuslean
\item The number $\thetaq$ is transcendental.  \statuslean
\item Every counterexample satisfies $x\ge12$ by the kernel-checked $4096$ barrier.
The separate interval computation in the combined report raises the checked barrier to
$x>20$ under its software trust assumptions.  \statuslean{} \statuscomp
\item The real four exponentials conjecture implies Conjecture~\ref{conj:AE}.
\statuslean
\end{itemize}

These facts eliminate rational, algebraic-irrational, small, and routine
three-exponentials cases.  They also explain why an elementary manipulation of logarithms is
unlikely to finish the proof: the unknown lies exactly where present transcendence theory
changes from theorem to conjecture.

\section{What failure would look like: the exact conditional normal form}
\label{sec:conditional}

The most important fact in the current repository is not another exclusion range.  It is the
complete description of all solutions if even one nonintegral solution exists.

\subsection{Canonical primitive generator}

Let
\[
 \Sol:=\{x\in\R:2^x,3^x\in\N\}.
\]
The theorem
\path{exists_canonical_primitiveGenerator_of_not_alaogluErdosConjecture}
proves the following conditional statement.

\begin{theorem}[Exact conditional solution monoid]
\label{thm:exact-monoid}
Assume Conjecture~\ref{conj:AE} is false.  Then there is a unique least nonintegral solution
$\beta>0$ such that
\begin{equation}
  \Sol=\{n+k\beta:n,k\in\Nzero\},
  \label{eq:solution-monoid}
\end{equation}
and the representation is unique.  If
\[
  M=2^\beta,\qquad A=3^\beta,
\]
then
\begin{equation}
  \Cand=\{2^nM^k:n,k\in\Nzero\},
  \label{eq:candidate-monoid}
\end{equation}
and this representation is also unique.
\end{theorem}

\begin{statusbox}{\statuslean}
This is a direct paraphrase of the unique-coordinate clauses in the cited primitive-generator
theorem.  The theorem also supplies the canonical odd core and radical normalization described
next.
\end{statusbox}

The phrase ``least nonintegral'' is stronger than choosing an arbitrary counterexample.  It
forces affine primitivity: no nontrivial common root can be extracted from the two output
integers while preserving the same base-adic offset.  This is the arithmetic analogue of
choosing a primitive lattice direction.

\subsection{Odd core, rational-power index, and localized radical}

The same theorem supplies natural numbers $w,d,a,c$ satisfying
\[
 \begin{gathered}
   w>1,\qquad w\text{ odd},\qquad d>0,\qquad c>0,\\
   \beta=a+d\log_2w,\\
   2^\beta=2^aw^d,\qquad 3^\beta=c,
 \end{gathered}
\]
with a gcd-normalization on the odd prime valuations of $w$.  Define
\[
  E:=w^{\thetaq}.
\]
Then
\begin{equation}
  E^d=\frac{c}{3^a},
  \qquad a=0\quad\text{or}\quad 3\nmid c,
  \label{eq:radical-normalization}
\end{equation}
and $d$ is the least positive exponent for which $E^d$ is rational.  The module
\path{RadicalDegree.lean} proves
\begin{equation}
 \operatorname{minpoly}_{\Q}(E)=X^d-\frac{c}{3^a},
 \qquad [\Q(E):\Q]=d,
 \label{eq:minpoly}
\end{equation}
with irreducibility of the displayed binomial.

Consequently the conjecture is kernel-verified equivalent to each of the following exclusion
statements.

\begin{enumerate}
\item For every odd $u>1$,
\begin{equation}
  u^{\thetaq}\notin\Z[1/3].
  \label{eq:localized-radical}
\end{equation}

\item It is enough to quantify over odd $u>1$ that are not proper powers.

\item There are no power-primitive odd $w>1$, $d>0$, and $a,B\in\Nzero$ such that
\begin{equation}
  3^a\bigl(w^{\thetaq}\bigr)^d=B.
  \label{eq:primitive-radical}
\end{equation}
\end{enumerate}

This is the cleanest remaining target.  It is only one real number at a time, but already the
special choice $u=3$ asks whether $3^{\log_2 3}$ can lie in $\Z[1/3]$, which is not decided by
current transcendence theorems.

\subsection{Simultaneous primitive normalization of both outputs}

The module \path{OutputNormalization.lean} gives a symmetric refinement.  For the least
$\beta$, one may write
\[
  M=2^aw^d,
  \qquad A=3^bv^e,
\]
where $w$ is odd and power-primitive, $3\nmid v$, $v$ is power-primitive, and
\begin{equation}
  a=0\quad\text{or}\quad b=0.
  \label{eq:one-depth-zero}
\end{equation}
Moreover,
\begin{equation}
 \gcd\!\left(d,e,|b-a|\right)=1.
 \label{eq:simultaneous-gcd}
\end{equation}
The outer exponents $d$ and $e$ are intrinsic: they are recovered as gcds of the prime
valuations of the base-coprime parts of $M$ and $A$.

These restrictions are strong, but they still admit infinitely many formal integer parameter
patterns.  The unresolved analytic equation remains
\[
  \frac{\log M}{\log2}=\frac{\log A}{\log3}.
\]
No known local valuation argument turns \eqref{eq:one-depth-zero} and
\eqref{eq:simultaneous-gcd} into a contradiction.

\subsection{Exact counts and the conditional population}

The monoid \eqref{eq:solution-monoid} makes the population of a hypothetical counterexample
completely explicit.  In every unit exponent block $[T,T+1)$,
\begin{equation}
 \#\bigl(\Sol\cap[T,T+1)\bigr)
 =1+\left\lfloor\frac{T+1}{\beta}\right\rfloor.
 \label{eq:unit-count}
\end{equation}
Passing through $m=2^x$, this becomes the exact dyadic magnitude count
\begin{equation}
 \#\bigl(\Cand\cap[2^T,2^{T+1})\bigr)
 =1+\left\lfloor\frac{T+1}{\beta}\right\rfloor.
 \label{eq:dyadic-count}
\end{equation}
Thus a failed conjecture would produce only linearly many candidates in the logarithmic block
index:
\[
 \#\bigl(\Cand\cap[X,2X]\bigr)\asymp\log X.
\]
Cumulatively,
\begin{equation}
 \#\bigl(\Cand\cap[1,2^T)\bigr)
 =\frac{T^2}{2\beta}+O(T).
 \label{eq:cumulative-quadratic}
\end{equation}
The repository formalizes both the exact floor sums and the limiting constant
$1/(2\beta)$.  It also proves the growth reformulation
\begin{equation}
 \text{Conjecture~\ref{conj:AE}}
 \iff
 \frac{\#(\Cand\cap[1,2^T))}{T^2}\longrightarrow0.
 \label{eq:growth-equivalence}
\end{equation}

This exact population is the correct benchmark for every proposed density or determinant
argument.  A method proving only $O((\log X)^2)$ candidates in $[X,2X]$, for example, is a
valid unconditional theorem but cannot contradict \eqref{eq:dyadic-count}.

\subsection{Blockwise irrational rotation}

For fixed $T$, the nonintegral solutions in $[T,T+1)$ have fractional parts drawn from the
initial orbit segment
\[
 \{\beta\},\{2\beta\},\ldots,
 \left\{\left\lfloor\frac{T+1}{\beta}\right\rfloor\beta\right\}.
\]
Because $\beta$ is irrational, the repository proves vanishing of every nonzero normalized
Fourier mode on these blocks.  The combined report derives the corresponding continuous-test
and interval-counting equidistribution at paper level.

This attractive regularity does not solve the conjecture.  The exact integrality condition is
not a fixed positive-measure target in the circle; it is an exponentially shrinking target.
That mismatch will recur below.

\section{The exact four-exponentials barrier}
\label{sec:four-exp}

\subsection{The $2\times2$ exponential matrix}

Suppose $x$ is a nonintegral solution.  Consider row vectors $(1,x)$ and column vectors
$(\log2,\log3)$.  The four corresponding exponentials are
\[
 \exp(\log2)=2,
 \quad \exp(\log3)=3,
 \quad \exp(x\log2)=2^x,
 \quad \exp(x\log3)=3^x.
\]
All four are algebraic, indeed positive integers.  The rows are linearly independent over
$\Q$ because $x$ is irrational; the columns are linearly independent because no nonzero
powers of $2$ and $3$ agree.  The real four exponentials conjecture predicts that at least one
of the four exponentials must be transcendental, a contradiction \cite{Waldschmidt2000,Waldschmidt2005}.

The Lean theorem
\path{alaogluErdosConjecture_of_realFourExponentialsConjecture}
formalizes exactly this implication.  It does not assume four exponentials as an axiom; it
defines the conjectural proposition and proves the conditional implication.

\subsection{Why six exponentials does not specialize}

The six exponentials theorem applies to either two independent row parameters and three
independent column parameters, or three rows and two columns.  Here only two canonical
logarithmic directions are available.  Adding a third base $5$ gives the theorem
\[
  2^x,3^x,5^x\in\Z\quad\Longrightarrow\quad x\in\Z,
\]
which is already formalized by an interpolation determinant.  In the two-base problem, any
additional base built from powers of $2$ and $3$ has logarithm in their rational span and does
not create the missing independent column.

The module \path{ThreeSmooth.lean} and the denominator-controlled results in
\path{RationalThirdBase.lean} make this obstruction precise: an auxiliary base with a prime
factor outside $\{2,3\}$ would be powerful, but its corresponding exponential value is not
forced by the original hypotheses.  Bases supported only on $2$ and $3$ do not increase the
transcendence rank.

\subsection{The integral special case may still be easier --- but not automatically}

Four exponentials asks only for algebraicity.  The present values are integers, so there is
extra discrete structure.  The determinant results in \cref{sec:arbitrary-node,sec:semigroup}
extract some of that structure: nonzero determinants are integers and therefore have absolute
value at least $1$.  However, the lower bound $1$ is the only universally available
archimedean information.  To beat the four-exponentials threshold, one likely needs to show
that the determinants are divisible by a large integer, or that a suitable common denominator
is anomalously small.  That observation motivates the $p$-adic research priority in
\cref{sec:future}.

\section{Full-proof attempts and their failure points}
\label{sec:attempts}

This section records the principal avenues explored before turning to the new determinant
results.  The purpose is not merely negative.  Each failure point identifies a quantitative or
structural feature that a successful proof must improve.

\subsection{Rational approximation to a least generator}

Assume failure and let $\beta$ be the least nonintegral solution, with
$M=2^\beta$ and $A=3^\beta$.  Let $p/q$ approximate $\beta$, and put
$\varepsilon=q\beta-p$.  Then
\begin{align}
 M^q-2^p
  &=2^p\left(e^{\varepsilon\log2}-1\right),
 \label{eq:approx-two}\\
 A^q-3^p
  &=3^p\left(e^{\varepsilon\log3}-1\right).
 \label{eq:approx-three}
\end{align}
Both left sides are integers.  If either is nonzero, its absolute value is at least $1$.
For small $\varepsilon$, this gives bounds of the shape
\[
 |\varepsilon|\gtrsim 2^{-p}
 \quad\text{and}\quad
 |\varepsilon|\gtrsim 3^{-p}.
\]
These are exponentially tiny in $q$.  Continued fractions only guarantee infinitely many
approximants with $|\varepsilon|\ll q^{-1}$, while Baker-type lower bounds for ordinary linear
forms in logarithms are polynomial in the relevant heights \cite{Baker1975,Waldschmidt2000}.  Neither scale conflicts with an
exponential lower bound as weak as $M^{-q}$ or $A^{-q}$.

One might hope to multiply \eqref{eq:approx-two} and \eqref{eq:approx-three}, or eliminate
$\varepsilon$ between them.  To first order their ratio merely recovers
$\log3/\log2$; higher terms are too small to create an integer contradiction.  The two
integer gaps are correlated, not independent.

\begin{observation}
Any rational-approximation proof needs more than ``a nonzero integer has size at least one.''
It needs a divisibility lower bound growing exponentially with $q$, or an approximation to
$\beta$ exponentially better than continued fractions can supply in general.
\end{observation}

\subsection{Equidistribution versus exponentially shrinking targets}

The multiples $\{k\beta\}$ are equidistributed modulo $1$.  Meanwhile, integrality of
$2^{n+t}$ between adjacent integral powers forces the fractional part $t$ away from both
endpoints by an amount comparable with $2^{-n}$; the base-$3$ condition gives a comparable
$3^{-n}$ restriction.  For the $k$th multiple, the relevant target width is therefore
exponentially small in $k$.

Equidistribution controls visits to fixed intervals and, with discrepancy estimates, to
moderately shrinking intervals.  It does not force visits to a summable sequence of targets
whose lengths behave like $e^{-ck}$.  Indeed, the Borel--Cantelli heuristic predicts only
finitely many such visits for a typical rotation.  The exact conditional monoid is perfectly
compatible with this behavior.

The repository's blockwise Weyl theorems are therefore structurally correct but point in the
wrong direction for a contradiction: they describe the distribution of already existing
solutions within each block, not the probability that arbitrary exponents hit the integer
lattice.

\subsection{Finite differences}

For $f(t)=t^{\thetaq}$, the repository proves the exact mean-value formula for every forward
difference order $r$:
\[
 \Delta_h^r f(n)
 =\fall{\thetaq}{r}h^r\xi^{\thetaq-r}
\qquad(n<\xi<n+rh).
\]
If all $r+1$ values $f(n),f(n+h),\ldots,f(n+rh)$ are integers, then the left side is an
integer.  Under
\[
 |\fall{\thetaq}{r}|h^r<n^{r-\thetaq},
\]
its absolute value is strictly between $0$ and $1$, a contradiction.  This is a sharp and
uniform obstruction to candidate arithmetic progressions.

The limitation is combinatorial.  A rank-two multiplicative monoid need not contain long
additive arithmetic progressions among its values.  Even many points in a short interval need
not be equally spaced.  The new arbitrary-node determinant in \cref{sec:arbitrary-node}
removes that equal-spacing hypothesis, but its packing exponent remains too weak to contradict
the exact conditional count.

\subsection{Radical fields and conjugates}

The normalized radical $E=w^{\thetaq}$ is algebraic and has exact minimal polynomial
$X^d-c/3^a$.  Its conjugates are $\zeta_d^jE$.  This gives complete algebraic control of the
field $\Q(E)$, its degree, and its norm.  What it does not give is a way to transport the real
analytic identity
\[
  \log E=\thetaq\log w
\]
to the other conjugates.  The chosen positive real logarithm is not Galois-equivariant; the
other roots differ by complex logarithmic periods.  Taking norms merely reproduces
$E^d=c/3^a$.

Ramification is similarly inconclusive.  The primes occurring in the binomial and in $w$ are
highly constrained, but the equation relates them through a quotient of logarithms, not
through an algebraic field embedding.  Standard Kummer theory sees the radical equation but
not the special value $\thetaq=\log3/\log2$.

\begin{observation}
A radical-field proof needs a genuinely mixed invariant: something simultaneously sensitive
to the algebraic conjugates of $E$ and to the logarithmic relation among $2,3,w,E$.  Separate
algebraic and real-analytic arguments currently meet only at the four-logarithm determinant.
\end{observation}

\subsection{Compact $S$-integral orbit}

From the least generator define
\[
 U_k:=\frac{M^k}{2^{\lfloor k\beta\rfloor}}
     =2^{\{k\beta\}}\in\Z[1/2],
 \qquad
 V_k:=\frac{A^k}{3^{\lfloor k\beta\rfloor}}
     =3^{\{k\beta\}}\in\Z[1/3].
 \label{eq:S-orbit}
\]
The points $(U_k,V_k)$ lie on the compact analytic arc
\[
 \mathcal A=\{(2^t,3^t):0\le t<1\}
\]
and are dense there because $\beta$ is irrational.  Their logarithmic heights grow linearly
in $k$, while the number of points of logarithmic height at most $H$ is also linear in $H$.

This suggests o-minimal and $S$-integral counting.  Butler's work relates special cases of
Wilkie-type algebraic-point counting to real three- and four-exponentials statements
\cite{Butler2017}.  The
standard Pila--Wilkie conclusion, however, is subpolynomial in the usual multiplicative height
outside algebraic parts.  Reparameterized by logarithmic height, the orbit's linear count is
not immediately contradictory, especially because the denominators are highly structured and
grow exponentially.

A promising refinement would count points on $\mathcal A$ with denominator supports exactly
$\{2\}$ and $\{3\}$ and with synchronized denominator exponents.  No theorem of the required
strength is currently available in the inspected corpus.

\subsection{Forcing an auxiliary base}

If one could exhibit an integer $b>1$ multiplicatively independent of $2$ and $3$ such that
$b^x$ is algebraic (or rational with suitably controlled denominator), then six exponentials
or the repository's rational-third-base theorems could force $x\in\Q$, hence $x\in\Z$.
The original hypotheses provide no such $b$.

The obvious constructions $b=2^u3^v$ remain in the same logarithmic span.  Constructions from
$M=2^x$ and $A=3^x$ also collapse to that span after taking logarithms.  A successful auxiliary
base must therefore arise from a nontrivial arithmetic operation whose exponential behavior is
controlled but whose logarithm leaves the rank-two span.  This is essentially an algebraic
independence problem in disguise.

\section{New result I: arbitrary-node lattice repulsion}
\label{sec:arbitrary-node}

The repository's higher-difference theorem uses equally spaced nodes.  The same integer-versus-
small-real principle extends to arbitrary nodes through interpolation determinants.

\subsection{The determinant identity}

For distinct real numbers $x_0<\cdots<x_r$ and a function $f$, write
$f[x_0,\ldots,x_r]$ for the divided difference.  The standard determinant identity is
\begin{equation}
 \det
 \begin{pmatrix}
  1&x_0&\cdots&x_0^{r-1}&f(x_0)\\
  1&x_1&\cdots&x_1^{r-1}&f(x_1)\\
  \vdots&\vdots&&\vdots&\vdots\\
  1&x_r&\cdots&x_r^{r-1}&f(x_r)
 \end{pmatrix}
 =V(x_0,\ldots,x_r)f[x_0,\ldots,x_r],
 \label{eq:det-divdiff}
\end{equation}
where
\[
 V(x_0,\ldots,x_r)=\prod_{0\le i<j\le r}(x_j-x_i).
\]
If $f\in C^r$ on the interval spanned by the nodes, the generalized mean-value theorem gives
some $\xi\in(x_0,x_r)$ such that
\begin{equation}
 f[x_0,\ldots,x_r]=\frac{f^{(r)}(\xi)}{r!}.
 \label{eq:divdiff-mvt}
\end{equation}
For $f(t)=t^{\thetaq}$,
\[
 f^{(r)}(t)=\fall{\thetaq}{r}t^{\thetaq-r},
 \qquad
 \fall{\thetaq}{r}:=\prod_{j=0}^{r-1}(\thetaq-j).
\]
Since $\thetaq$ is not an integer, the coefficient never vanishes.

\subsection{Arbitrary-node theorem}

\begin{theorem}[Vandermonde repulsion for candidates]
\label{thm:arbitrary-node}
Let $r\ge2$ and let
\[
 1\le m_0<m_1<\cdots<m_r
\]
be natural candidates, so $m_i^{\thetaq}\in\N$.  Then
\begin{equation}
 \prod_{0\le i<j\le r}(m_j-m_i)
 \ge
 \frac{r!}{|\fall{\thetaq}{r}|}\,m_0^{r-\thetaq}.
 \label{eq:arbitrary-node-bound}
\end{equation}
\end{theorem}

\begin{proof}
Apply \eqref{eq:det-divdiff} to $f(t)=t^{\thetaq}$ and the nodes $m_i$.  Every entry of the
matrix is an integer: the ordinary powers are integral by construction and the last column is
integral because the $m_i$ are candidates.  By \eqref{eq:divdiff-mvt}, its determinant is
\[
 D=V(m_0,\ldots,m_r)
   \frac{\fall{\thetaq}{r}}{r!}\xi^{\thetaq-r}
\qquad(m_0<\xi<m_r).
\]
It is nonzero, so $D\in\Z\setminus\{0\}$ and $|D|\ge1$.  Because $r\ge2>\thetaq$, the exponent
$\thetaq-r$ is negative and
\[
 \xi^{\thetaq-r}\le m_0^{\thetaq-r}.
\]
Therefore
\[
 1\le |D|
 \le V(m_0,\ldots,m_r)
      \frac{|\fall{\thetaq}{r}|}{r!}m_0^{\thetaq-r},
\]
which rearranges to \eqref{eq:arbitrary-node-bound}.
\end{proof}

\begin{statusbox}{\statusnew}
The argument is unconditional and elementary once the divided-difference mean-value theorem
is available.  It is not present in the inspected Lean modules; \path{HigherDifference.lean}
uses arithmetic progressions.  A direct formalization blueprint appears in
\cref{sec:lean-blueprint}.
\end{statusbox}

\subsection{Three points: a lattice-triangle inequality}

For $r=2$, the determinant is
\[
 D=\det\begin{pmatrix}
 1&m_0&m_0^{\thetaq}\\
 1&m_1&m_1^{\thetaq}\\
 1&m_2&m_2^{\thetaq}
 \end{pmatrix}.
\]
Geometrically, $|D|$ is twice the area of the lattice triangle with vertices
$(m_i,m_i^{\thetaq})$.  Strict convexity makes the area positive; lattice integrality makes it
at least $1/2$.

\begin{corollary}[Three-candidate span]
\label{cor:three-span}
If three candidates lie in $[N,N+H]$, where $N\ge1$, then
\begin{equation}
 H\ge
 \left(\frac{8}{\thetaq(\thetaq-1)}\right)^{1/3}
 N^{(2-\thetaq)/3}.
 \label{eq:three-span}
\end{equation}
Numerically,
\[
 H\ge2.0510723956\ldots\,N^{0.1383458330\ldots}.
\]
\end{corollary}

\begin{proof}
Write $a=m_1-m_0$, $b=m_2-m_1$.  Then $a,b>0$, $a+b\le H$, and
\[
 V=a b(a+b)\le\frac{H^3}{4},
\]
with the continuous maximum at $a=b=H/2$.  Theorem~\ref{thm:arbitrary-node} at $r=2$ gives
\[
 V\ge\frac{2}{\thetaq(\thetaq-1)}N^{2-\thetaq}.
\]
Combining the inequalities yields \eqref{eq:three-span}.
\end{proof}

\subsection{General local packing bound}

Put
\[
 R_r:=\binom{r+1}{2}=\frac{r(r+1)}2,
 \qquad
 C_r:=\left(\frac{r!}{|\fall{\thetaq}{r}|}\right)^{1/R_r},
 \qquad
 \alpha_r:=\frac{r-\thetaq}{R_r}.
\]
Since every pairwise gap among points in $[N,N+H]$ is at most $H$, the Vandermonde product is
at most $H^{R_r}$.  Theorem~\ref{thm:arbitrary-node} immediately gives:

\begin{corollary}[At most $r$ points below the repulsion scale]
\label{cor:r-packing}
If $r+1$ candidates lie in $[N,N+H]$, then
\begin{equation}
 H\ge C_rN^{\alpha_r}.
 \label{eq:r-span}
\end{equation}
Equivalently, every interval of length strictly less than $C_rN^{\alpha_r}$ beginning at or
after $N$ contains at most $r$ candidates.
\end{corollary}

The exponent can be optimized.  As a function of a real variable $r>\thetaq$,
\[
 \alpha(r)=\frac{2(r-\thetaq)}{r(r+1)}
\]
has its maximum at
\[
 r=\thetaq+\sqrt{\thetaq^2+\thetaq}=3.60908\ldots.
\]
Hence the best integer order in this ordinary-power family is $r=4$.

\begin{table}[ht]
\centering
\small
\begin{tabular}{rrrrr}
\toprule
$r$ & $R_r$ & $|\fall{\thetaq}{r}|$ & $\alpha_r$ & $C_r$ \\
\midrule
2 & 3  & 0.9271436280 & 0.1383458331 & 1.2920946431 \\
3 & 6  & 0.3847993728 & 0.2358395832 & 1.5805909191 \\
4 & 10 & 0.5445055422 & 0.2415037499 & 1.4602287461 \\
5 & 15 & 1.3150013031 & 0.2276691666 & 1.3510881062 \\
6 & 21 & 4.4907787617 & 0.2102398809 & 1.2735045942 \\
7 & 28 & 19.8269566337 & 0.1933941964 & 1.2187063909 \\
8 & 36 & 107.3637136680 & 0.1781954861 & 1.1790125192 \\
\bottomrule
\end{tabular}
\caption{Ordinary-power arbitrary-node repulsion constants.  For $r=2$, the sharper
three-point constant in \cref{cor:three-span} uses the exact maximum of the three-gap product.}
\label{tab:ordinary-constants}
\end{table}

In particular:

\begin{corollary}[Best ordinary-power local statement]
\label{cor:four-span}
Every interval
\[
 [N,N+H],
 \qquad H<1.4602287460\,N^{0.2415037499},
\]
contains at most four natural candidates.
\end{corollary}

\subsection{A bound for any longer interval}

The span estimate can be converted into a counting inequality.  Let
$m_1<\cdots<m_s$ be all candidates in $[N,N+H]$.  For each $i\le s-r$, the consecutive
window $m_i,\ldots,m_{i+r}$ has span at least $C_rN^{\alpha_r}$.  Summing these spans, each
adjacent gap $m_{j+1}-m_j$ is counted at most $r$ times.  Hence
\[
 (s-r)C_rN^{\alpha_r}
 \le r(m_s-m_1)
 \le rH.
\]

\begin{corollary}[Arbitrary-node counting inequality]
\label{cor:ordinary-count}
For every $r\ge2$,
\begin{equation}
 \#(\Cand\cap[N,N+H])
 \le r+\frac{r}{C_r}H N^{-\alpha_r}.
 \label{eq:ordinary-count}
\end{equation}
\end{corollary}

This is a genuine strengthening of progression avoidance: it constrains every possible
configuration.  It is still far from the conditional benchmark.  At the optimal $r=4$ and
$H=N$, \eqref{eq:ordinary-count} is $O(N^{0.7585\ldots})$, whereas the exact failed-conjecture
population would be only $O(\!\log N)$.

\subsection{Relation to higher differences}

For equally spaced nodes $m_i=n+ih$, the Vandermonde factor is explicit:
\[
 V=h^{R_r}\prod_{j=1}^{r}j!.
\]
Substitution in \eqref{eq:det-divdiff} recovers an integer interpolation remainder closely
related to the $r$th forward difference.  The repository theorem is sharper for this special
configuration because it uses the exact binomial combination rather than bounding all gaps
by a common span.  Conversely, Theorem~\ref{thm:arbitrary-node} applies when the nodes have no
additive structure at all.  The two results are complementary.

\section{New result II: semigroup generalized Vandermonde determinants}
\label{sec:semigroup}

The preceding determinant used only the integral columns
$1,m,m^2,\ldots$ and one extra integral column $m^{\thetaq}$.  But a candidate carries much
more integral data.  If $m^{\thetaq}=n\in\N$, then for every $a,b\in\Nzero$,
\begin{equation}
 m^{a+b\thetaq}=m^a n^b\in\N.
 \label{eq:mixed-integrality}
\end{equation}
The full set of available exponents is therefore
\[
 \LambdaTheta:=\{a+b\thetaq:a,b\in\Nzero\}.
\]
Since $\thetaq$ is irrational, all pairs $(a,b)$ give distinct exponents.

\subsection{Positivity of generalized Vandermonde determinants}

Let
\[
 0\le\lambda_0<\lambda_1<\cdots<\lambda_r
\qquad\text{and}\quad
 0<x_0<x_1<\cdots<x_r.
\]
The functions $x^{\lambda_j}$ form an extended complete Chebyshev system on $(0,\infty)$.
Their Wronskian is \cite{KarlinStudden}
\begin{equation}
 W(x)
 =x^{\sum_j\lambda_j-R_r}
   \prod_{0\le i<j\le r}(\lambda_j-\lambda_i)>0.
 \label{eq:wronskian}
\end{equation}
Indeed, after factoring $x^{\lambda_j}$ from column $j$ and $x^{-i}$ from derivative row
$i$, the remaining determinant is the Vandermonde determinant of the $\lambda_j$, expressed
in the falling-factorial basis.  Positivity of all initial Wronskians implies strict positivity
of every ordered evaluation determinant.

For completeness, there is also an elementary zero-count proof.  Put $x=e^t$.  A nonzero
linear combination of $x^{\lambda_j}$ becomes an exponential polynomial
$\sum c_je^{\lambda_jt}$.  After factoring the lowest exponential, differentiation removes
the constant term and leaves one fewer exponential.  Induction with Rolle's theorem shows
that a combination of $r+1$ such functions has at most $r$ real zeros counting multiplicity.
Consequently an evaluation determinant at $r+1$ ordered nodes cannot vanish.  Its sign is
constant on the connected chamber $x_0<\cdots<x_r$ and is positive in the coalescing-node
limit by \eqref{eq:wronskian}.

\begin{lemma}[Generalized Vandermonde positivity]
\label{lem:generalized-vandermonde}
For strictly increasing nonnegative exponents $\lambda_j$ and strictly increasing positive
nodes $x_i$,
\begin{equation}
 \det\bigl(x_i^{\lambda_j}\bigr)_{0\le i,j\le r}>0.
 \label{eq:gv-positive}
\end{equation}
\end{lemma}

\subsection{Integer determinants from candidates}

\begin{theorem}[Semigroup integer determinant]
\label{thm:semigroup-integer-det}
Let $m_0<\cdots<m_r$ be natural candidates.  Choose distinct pairs
$(a_j,b_j)\in\Nzero^2$ and order
\[
 \lambda_j=a_j+b_j\thetaq
\]
so that $\lambda_0<\cdots<\lambda_r$.  Then
\begin{equation}
 D_{\boldsymbol\lambda}(\boldsymbol m)
 :=\det\bigl(m_i^{\lambda_j}\bigr)_{0\le i,j\le r}
 \label{eq:semigroup-det}
\end{equation}
is a positive integer.  In particular, $D_{\boldsymbol\lambda}(\boldsymbol m)\ge1$.
\end{theorem}

\begin{proof}
By \eqref{eq:mixed-integrality}, each entry equals
\[
 m_i^{a_j}(m_i^{\thetaq})^{b_j}\in\N,
\]
so the determinant is an integer.  It is strictly positive by
Lemma~\ref{lem:generalized-vandermonde}.
\end{proof}

\begin{statusbox}{\statusnew}
This theorem packages all mixed integrality consequences into one reusable object.  It is
not the same as the repository's exponential interpolation determinant in
\path{MultiplicativeRank.lean}: here the rows are arbitrary natural candidates and the
columns are chosen from the two-dimensional exponent semigroup.
\end{statusbox}

\subsection{Divided-difference factorization}

Let $f_j(x)=x^{\lambda_j}$ and let
\[
 V(\boldsymbol m)=\prod_{i<j}(m_j-m_i).
\]
Applying the Newton divided-difference row transformation gives
\begin{equation}
 D_{\boldsymbol\lambda}(\boldsymbol m)
 =V(\boldsymbol m)
  \det\bigl(f_j[m_0,\ldots,m_i]\bigr)_{0\le i,j\le r}.
 \label{eq:newton-factorization}
\end{equation}
The transformation is exact; its determinant is $V(\boldsymbol m)^{-1}$.

Suppose
\[
 N\le m_0<\cdots<m_r\le N+H\le2N.
\]
For each divided difference, the mean-value theorem gives a point in $[N,2N]$ and
\begin{equation}
 |f_j[m_0,\ldots,m_i]|
 \le
 \frac{|\fall{\lambda_j}{i}|}{i!}
 \sup_{N\le x\le2N}x^{\lambda_j-i}.
 \label{eq:dd-first-bound}
\end{equation}
We need a uniform elementary coefficient estimate.

\begin{lemma}[Falling-factorial coefficient bound]
\label{lem:falling-bound}
For $\lambda\ge0$ and $i\in\Nzero$,
\begin{equation}
 \frac{|\fall{\lambda}{i}|}{i!}\le2^{\lambda+i}.
 \label{eq:falling-bound}
\end{equation}
\end{lemma}

\begin{proof}
If $0\le\lambda\le1$, each factor satisfies
$|\lambda-k|\le k+1$, so the quotient is at most $1$.
If $\lambda\ge1$, let $n=\lceil\lambda\rceil$.  Then
\[
 \frac{|\fall{\lambda}{i}|}{i!}
 \le\prod_{k=0}^{i-1}\frac{\lambda+k}{k+1}
 \le\prod_{k=0}^{i-1}\frac{n+k}{k+1}
 =\binom{n+i-1}{i}.
\]
The last binomial coefficient is at most the sum of its row,
$2^{n+i-1}$, and $n-1\le\lambda$.  Hence it is at most $2^{\lambda+i}$.
\end{proof}

For $x\in[N,2N]$,
\begin{equation}
 x^{\lambda-i}\le2^\lambda N^{\lambda-i}.
 \label{eq:power-strip-bound}
\end{equation}
If $\lambda-i\ge0$, use $x\le2N$; if it is negative, use $x\ge N$.
Combining \eqref{eq:dd-first-bound}, \eqref{eq:falling-bound}, and
\eqref{eq:power-strip-bound},
\begin{equation}
 |f_j[m_0,\ldots,m_i]|
 \le2^{2\lambda_j+i}N^{\lambda_j-i}.
 \label{eq:dd-final-bound}
\end{equation}

\subsection{Quantitative semigroup gap theorem}

Set
\[
 R:=R_r=\frac{r(r+1)}2,
 \qquad
 \Sigma:=\sum_{j=0}^{r}\lambda_j.
\]
By the Leibniz formula and \eqref{eq:dd-final-bound}, every permutation term in the small
determinant in \eqref{eq:newton-factorization} is at most
\[
 2^{2\Sigma+R}N^{\Sigma-R},
\]
and there are $(r+1)!$ terms.  Since every pairwise node difference is at most $H$,
$V(\boldsymbol m)\le H^R$.  The positive integer lower bound from
Theorem~\ref{thm:semigroup-integer-det} now gives the following.

\begin{theorem}[Semigroup determinant gap]
\label{thm:semigroup-gap}
Let $r\ge1$.  Let $\lambda_0<\cdots<\lambda_r$ be any $r+1$ exponents in
$\LambdaTheta$, and put $R=r(r+1)/2$ and $\Sigma=\sum_j\lambda_j$.  If $r+1$ natural
candidates lie in
\[
 [N,N+H]\subseteq[N,2N],
 \qquad N\ge1,
\]
then
\begin{equation}
 H\ge
 ((r+1)!)^{-1/R}
 2^{-1-2\Sigma/R}
 N^{1-\Sigma/R}.
 \label{eq:semigroup-gap}
\end{equation}
\end{theorem}

\begin{proof}
Equations \eqref{eq:newton-factorization} and \eqref{eq:dd-final-bound} imply
\[
 1\le D_{\boldsymbol\lambda}(\boldsymbol m)
 \le
 H^R(r+1)!2^{R+2\Sigma}N^{\Sigma-R}.
\]
Taking $R$th roots and rearranging proves \eqref{eq:semigroup-gap}.
\end{proof}

The exponent $1-\Sigma/R$ is the key.  We should choose the $r+1$ smallest available
semigroup exponents to minimize $\Sigma$.

\subsection{The ordered exponent semigroup}

Write the elements of $\LambdaTheta$ in increasing order:
\[
 0=\lambda_0<\lambda_1<\lambda_2<\cdots.
\]
The beginning is
\begin{align*}
 0,\ 1,\ \thetaq,\ 2,\ 1+\thetaq,\ 3,\ 2\thetaq,\ 2+\thetaq,\ 4,
 \ 1+2\thetaq,\ 3+\thetaq,\ 3\thetaq,\ 5,\ldots
\end{align*}
The counting function is elementary:
\begin{align}
 L(T)
 &:=\#\{\lambda\in\LambdaTheta:\lambda\le T\}\notag\\
 &=\sum_{0\le b\le T/\thetaq}
   \left(\left\lfloor T-b\thetaq\right\rfloor+1\right)
 =\frac{T^2}{2\thetaq}+O(T).
 \label{eq:semigroup-counting}
\end{align}
Irrationality of $\thetaq$ ensures that different pairs $(a,b)$ are counted distinctly.
Inverting \eqref{eq:semigroup-counting} and summing the order statistics gives
\begin{align}
 \lambda_r
 &=\sqrt{2\thetaq r}+O(1),
 \label{eq:lambda-asymp}\\
 \Sigma_r:=\sum_{j=0}^{r}\lambda_j
 &=\frac23\sqrt{2\thetaq}\,r^{3/2}+O(r),
 \label{eq:sigma-asymp}\\
 \frac{\Sigma_r}{R_r}
 &=\frac43\sqrt{\frac{2\thetaq}{r}}+O(r^{-1}).
 \label{eq:sigma-ratio-asymp}
\end{align}
Consequently the power of $N$ in \eqref{eq:semigroup-gap} tends to $1$:
\begin{equation}
 1-\frac{\Sigma_r}{R_r}
 =1-\frac43\sqrt{\frac{2\thetaq}{r}}+O(r^{-1}).
 \label{eq:delta-asymp}
\end{equation}

\begin{table}[ht]
\centering
\small
\begin{tabular}{rcc}
\toprule
$r$ & $\Sigma_r$ & $1-\Sigma_r/R_r$ \\
\midrule
2  & $2.5849625007$  & $0.1383458331$ \\
3  & $4.5849625007$  & $0.2358395832$ \\
4  & $7.1699250014$  & $0.2830074999$ \\
5  & $10.1699250014$ & $0.3220049999$ \\
6  & $13.3398500029$ & $0.3647690475$ \\
7  & $16.9248125036$ & $0.3955424106$ \\
8  & $20.9248125036$ & $0.4187552082$ \\
9  & $25.0947375050$ & $0.4423391666$ \\
10 & $29.6797000058$ & $0.4603690908$ \\
12 & $39.4345875079$ & $0.4944283653$ \\
20 & $87.7939875195$ & $0.5819333928$ \\
\bottomrule
\end{tabular}
\caption{Exponent gain from using the first $r+1$ mixed exponents in
$\Nzero+\Nzero\thetaq$.  The ordinary-power family peaks and then decays; the semigroup
family improves monotonically in these data and tends to exponent $1$.}
\label{tab:semigroup-delta}
\end{table}

\subsection{An explicit coarse form}

For the global counting corollary we need no delicate lattice-point error term.  A simple
bound suffices.

\begin{lemma}[Crude semigroup weight bound]
\label{lem:crude-semigroup}
For $r\ge2$, the first $r+1$ semigroup exponents satisfy
\begin{equation}
 \lambda_r\le6\sqrt r,
 \qquad
 \frac{\Sigma_r}{R_r}\le\frac{12}{\sqrt r}.
 \label{eq:crude-semigroup}
\end{equation}
\end{lemma}

\begin{proof}
Let $q=\lceil\sqrt{r+1}\rceil$.  The $(q+1)^2$ pairs
$0\le a,b\le q$ give distinct exponents and each is at most
$q(1+\thetaq)<3q$.  Since $(q+1)^2>r+1$, the $r$th ordered exponent is below $3q$.
For $r\ge2$,
$q\le\sqrt{r+1}+1\le2\sqrt r$, hence $\lambda_r\le6\sqrt r$.
Then
\[
 \Sigma_r\le(r+1)\lambda_r\le6(r+1)\sqrt r,
\]
and division by $R_r=r(r+1)/2$ gives the second inequality.
\end{proof}

For $r\ge144$, \eqref{eq:crude-semigroup} gives $\Sigma_r/R_r\le1$.  Also
\[
 ((r+1)!)^{1/R_r}
 \le(r+1)^{(r+1)/R_r}
 =(r+1)^{2/r}\le3
\]
for $r\ge2$.  Therefore Theorem~\ref{thm:semigroup-gap} yields a particularly clean form.

\begin{corollary}[Explicit near-linear span]
\label{cor:explicit-semigroup-span}
If $r\ge144$ and $r+1$ natural candidates lie in
$[N,N+H]\subseteq[N,2N]$, then
\begin{equation}
 H\ge\frac1{24}N^{1-12/\sqrt r}.
 \label{eq:explicit-semigroup-span}
\end{equation}
\end{corollary}

\begin{proof}
In \eqref{eq:semigroup-gap}, use
$((r+1)!)^{-1/R}\ge1/3$,
$2^{-1-2\Sigma/R}\ge2^{-3}=1/8$, and, since $N\ge1$,
$N^{1-\Sigma/R}\ge N^{1-12/\sqrt r}$.
\end{proof}

\subsection{Counting candidates in a general interval}

As in \cref{cor:ordinary-count}, sum the spans of consecutive windows of $r+1$ candidates.
Every adjacent gap is counted at most $r$ times.

\begin{corollary}[Semigroup interval count]
\label{cor:semigroup-interval-count}
For $r\ge144$ and $1\le H\le N$,
\begin{equation}
 \#(\Cand\cap[N,N+H])
 \le r+24rH N^{-1+12/\sqrt r}.
 \label{eq:semigroup-interval-count}
\end{equation}
\end{corollary}

\begin{proof}
If $s$ is the number of candidates and $s\le r$, there is nothing to prove.  Otherwise,
Corollary~\ref{cor:explicit-semigroup-span} gives a lower bound
$L=N^{1-12/\sqrt r}/24$ for each span $m_{i+r}-m_i$.  Thus
\[
 (s-r)L\le rH,
\]
which is equivalent to \eqref{eq:semigroup-interval-count}.
\end{proof}

\subsection{An independent dyadic logarithmic-square bound}

\begin{corollary}[Explicit $O((\log X)^2)$ dyadic bound]
\label{cor:semigroup-log-square}
For every real $X\ge1$,
\begin{equation}
 \#(\Cand\cap[X,2X])
 \le10000\max\{1,\log X\}^2.
 \label{eq:semigroup-log-square}
\end{equation}
\end{corollary}

\begin{proof}
Put $L=\max\{1,\log X\}$ and
\[
 r=\left\lceil144L^2\right\rceil.
\]
Then $r\ge144$, $\sqrt r\ge12L$, and
\[
 X^{12/\sqrt r}
 =\exp\!\left(\frac{12\log X}{\sqrt r}\right)
 \le e.
\]
Apply \eqref{eq:semigroup-interval-count} with $N=H=X$:
\[
 \#(\Cand\cap[X,2X])
 \le r+24rX^{12/\sqrt r}
 \le(1+24e)r.
\]
Since $L\ge1$, $r\le145L^2$, and
$(1+24e)145<10000$.
\end{proof}

\begin{statusbox}{\statusnew}
The bound is numerically coarse by design.  Its value is conceptual independence: it comes
from positivity and integrality of generalized Vandermonde determinants, not from the
prime-factor valuation rank of candidates.  Refining constants is much less important than
finding an additional source of determinant lower bound.
\end{statusbox}

\section{Why the new hierarchy still stops short}
\label{sec:deficit}

The semigroup determinant is much stronger locally than the ordinary-power determinant: its
repulsion exponent approaches $1$.  It is therefore important to identify exactly why it does
not prove the conjecture.

\subsection{How many nodes are needed to approach scale $X$?}

The asymptotic exponent in \eqref{eq:delta-asymp} is
\[
 \delta_r:=1-\frac{\Sigma_r}{R_r}
 =1-\frac{c_\theta}{\sqrt r}+O(r^{-1}),
 \qquad
 c_\theta=\frac43\sqrt{2\thetaq}=2.373\ldots.
\]
For nodes near $X$, the lower span has the shape
\[
 X^{\delta_r}
 =X\exp\!\left(-\frac{c_\theta\log X}{\sqrt r}+O\!\left(\frac{\log X}{r}\right)\right).
 \label{eq:scale-loss}
\]
To make the exponential loss in \eqref{eq:scale-loss} bounded independently of $X$, one needs
\[
 \sqrt r\asymp\log X,
 \qquad\text{that is,}\qquad
 r\asymp(\log X)^2.
 \label{eq:node-threshold}
\]
This is precisely why Corollary~\ref{cor:semigroup-log-square} has a logarithmic square.

\subsection{How many nodes a counterexample supplies}

Conditional on a counterexample, the exact block count \eqref{eq:dyadic-count} gives
\[
 \#(\Cand\cap[X,2X])
 =\frac{\log_2X}{\beta}+O(1)
 \asymp\log X.
 \label{eq:conditional-log-population}
\]
Thus the actual conditional population is smaller than the determinant threshold
\eqref{eq:node-threshold} by one factor of $\log X$.

Taking $r\asymp\log X$ in \eqref{eq:scale-loss} yields only
\[
 X\exp(-C\sqrt{\log X}),
\]
which is much smaller than the average spacing $X/\log X$ among the
$\asymp\log X$ conditional candidates.  The inequality is therefore compatible with the
exact monoid even at the level of expected spacing.

\subsection{Broader multiplicative intervals do not help}

The cumulative conditional count below $X$ is $\asymp(\log X)^2$, so one might try to use all
those points at once.  But their nodes range from $1$ to $X$, and the gap theorem depends on
the smallest node.  Restricting to $[X^\gamma,X]$ with fixed $0<\gamma<1$ does retain
$\asymp(\log X)^2$ points, but its additive length is $\asymp X$, while the lower scale is at
most $X^{\gamma+o(1)}$.  Again there is no contradiction.

This is not a mere defect of coarse constants.  It reflects a mismatch between additive node
geometry and the logarithmic triangular lattice of exponents $(n,k)$ in
$m=2^nM^k$.

\subsection{The determinant needs one more source of size}

The archimedean argument uses only
\[
 D\in\Z_{>0}\quad\Longrightarrow\quad D\ge1.
\]
If the exact monoid coordinates forced
\[
 D\equiv0\pmod{Q}
\qquad\text{with}\quad
 \log Q\gg R\log X/\sqrt r,
\]
then the lower bound $D\ge Q$ could remove the one-logarithm deficit.  Alternatively, one could
seek a determinant with total exponent weight $\Sigma=O(r)$ rather than
$\Sigma\asymp r^{3/2}$, but the two-dimensional semigroup counting law
$L(T)\asymp T^2$ makes that impossible using only distinct nonnegative mixed monomials.

This gives a concrete design specification:

\begin{quote}
A successful determinant proof must add nonarchimedean divisibility, introduce lower-weight
functions outside the mixed-monomial semigroup, or exploit the exact two-dimensional row grid
so efficiently that $O(\!\log X)$ nodes suffice.
\end{quote}

\section{Further equivalent reformulations}
\label{sec:reformulations}

Several equivalent forms are already formalized in the repository.  The exact conditional
structure also yields useful new paper-level reformulations that expose different possible
proof technologies.

\subsection{Integer points on a transcendental lattice curve}

The candidate formulation \eqref{eq:candidate-equivalence} is
\begin{equation}
 \{(m,n)\in\N^2:n=m^{\thetaq}\}
 =\{(2^k,3^k):k\in\Nzero\}.
 \label{eq:lattice-curve}
\end{equation}
The curve is strictly convex and transcendental.  Theorem~\ref{thm:arbitrary-node} can be read
as a lower bound for lattice simplex volumes cut from this curve.

\subsection{Vanishing determinant of four logarithms}

As in \eqref{eq:log-det}, the conjecture is equivalent to
\begin{equation}
 \log M\log3-\log A\log2=0,
 \quad M,A\in\N,
 \quad\Longrightarrow\quad
 (M,A)=(2^k,3^k).
 \label{eq:four-log-form}
\end{equation}
This is the most direct bridge to four exponentials and algebraic independence of logarithms.

\subsection{Primitive localized-radical exclusion}

The kernel-verified radical form is
\begin{equation}
 \forall\,w>1\text{ odd and power-primitive},\quad
 \forall\,d>0,\ a\ge0,\qquad
 3^a(w^{\thetaq})^d\notin\N.
 \label{eq:primitive-localized}
\end{equation}
If this fails with least rational-power index $d$, then
$X^d-c/3^a$ is the exact minimal polynomial of $w^{\thetaq}$.

This formulation is ideal for Kummer, norm, and ramification investigations because all
redundant proper powers have been removed and the algebraic degree is known exactly.

\subsection{Exact growth dichotomy}

The repository already formalizes the cumulative subquadratic equivalence
\eqref{eq:growth-equivalence}.  The exact block formula gives several immediate variants.

\begin{proposition}[Block-growth reformulations]
\label{prop:block-reforms}
Conjecture~\ref{conj:AE} is equivalent to each of the following:
\begin{enumerate}
\item the dyadic block counts
$B(T)=\#(\Cand\cap[2^T,2^{T+1}))$ are bounded;
\item $B(T)=o(T)$;
\item $\liminf_{T\to\infty}B(T)/T=0$;
\item $B(T)=1$ for every $T\in\Nzero$.
\end{enumerate}
\end{proposition}

\begin{proof}
If the conjecture holds, every block contains only $2^T$, so all four properties hold.  If it
fails, \eqref{eq:dyadic-count} gives
$B(T)=1+\lfloor(T+1)/\beta\rfloor$, so $B(T)/T\to1/\beta>0$ and none of the first three
properties holds.  The fourth is visibly equivalent to the candidate statement.
\end{proof}

This is useful because a future analytic theorem need not classify candidates individually.
Any uniform $o(T)$ bound in the $T$th dyadic block would prove the conjecture.

\subsection{Compact synchronized $S$-integral orbit}

The orbit \eqref{eq:S-orbit} yields another equivalence.

\begin{proposition}[Compact-orbit reformulation]
\label{prop:compact-orbit}
The conjecture fails if and only if there exist integers $M,A>1$ and an irrational
$\beta>0$ such that, for every $k\ge1$, the point
\begin{equation}
 \left(
  \frac{M^k}{2^{\lfloor k\beta\rfloor}},
  \frac{A^k}{3^{\lfloor k\beta\rfloor}}
 \right)
 \label{eq:compact-orbit-points}
\end{equation}
lies on the compact arc
\[
 \mathcal A=\{(2^t,3^t):0\le t<1\}.
\]
Under these conditions the orbit is dense in $\mathcal A$.
\end{proposition}

\begin{proof}
Failure supplies the least generator $\beta$, with $M=2^\beta$ and $A=3^\beta$; then
\eqref{eq:compact-orbit-points} is exactly $(2^{\{k\beta\}},3^{\{k\beta\}})$, and density
follows from irrational rotation.

Conversely, arc membership gives $t_k\in[0,1)$ such that
\[
 k\log_2M-\lfloor k\beta\rfloor=t_k
 =k\log_3A-\lfloor k\beta\rfloor.
\]
Thus $\log_2M=\log_3A=:\gamma$, and
$0\le k\gamma-\lfloor k\beta\rfloor<1$ for every $k$.  Hence
$\lfloor k\gamma\rfloor=\lfloor k\beta\rfloor$, so
$|k(\gamma-\beta)|<1$ for every $k$; therefore $\gamma=\beta$.  Consequently
$M=2^\beta$ and $A=3^\beta$ are integers while $\beta$ is nonintegral.  The same identity now
shows that the orbit equals the irrational-rotation orbit and is dense.
\end{proof}

This formulation isolates the synchronized-denominator feature that a specialized
$S$-integral counting theorem would need to exploit.

\subsection{Two-dimensional exponential grid}

Under failure, every solution coordinate $(n,k)$ gives
\[
 m_{n,k}=2^nM^k,
 \qquad
 m_{n,k}^{\thetaq}=3^nA^k.
\]
For every mixed exponent $a+b\thetaq$,
\begin{equation}
 m_{n,k}^{a+b\thetaq}
 =(2^a3^b)^n(M^aA^b)^k.
 \label{eq:two-dimensional-grid}
\end{equation}
Thus the candidate evaluation matrix factors through a two-dimensional row lattice and a
two-dimensional column lattice.  The conjecture can be viewed as the assertion that this
integer exponential grid cannot exist unless $(M,A)$ is a common integral power of $(2,3)$.
This is the most natural starting point for a bespoke $2\times2$ interpolation determinant
that uses more than arbitrary-node geometry.

\section{Promising future research directions}
\label{sec:future}

The priorities below are ordered by how directly they address the quantitative deficit in
\cref{sec:deficit}.

\subsection{Priority 1: $p$-adic amplification of semigroup determinants}

The determinant
\[
 D=\det\bigl(m_i^{a_j}(m_i^{\thetaq})^{b_j}\bigr)
\]
is a positive integer.  Under the exact monoid,
\[
 m_i=2^{n_i}M^{k_i},
 \qquad
 m_i^{\thetaq}=3^{n_i}A^{k_i}.
\]
Every entry therefore has explicit valuations
\begin{align*}
 \vpp{2}{\text{entry}_{ij}}
 &=a_j(n_i+\vpp{2}{M}k_i)+b_j\vpp{2}{A}k_i,\\
 \vpp{3}{\text{entry}_{ij}}
 &=a_j\vpp{3}{M}k_i+b_j(n_i+\vpp{3}{A}k_i).
\end{align*}
A naive row or column factor is lost because the exponent set contains $(0,0)$.  More subtle
min-plus determinant estimates may nevertheless force divisibility after suitable row/column
finite-difference operations.  The exact affine primitivity and the condition
$\gcd(d,e,|b-a|)=1$ should be incorporated rather than discarded.

A concrete target is:
\begin{question}
Can one choose $O(\!\log X)$ candidate rows and mixed-monomial columns so that the resulting
nonzero determinant has
\[
 \log|D|_p^{-1}\gg(\log X)^2
\]
for $p=2$ or $3$, while the archimedean estimate is
$\log|D|_\infty=o((\log X)^2)$?
\end{question}
A positive answer would activate the product formula or simply the divisibility lower bound
$|D|\ge p^L$ and could erase the missing logarithm.

\subsection{Priority 2: exploit the full $(n,k)$ grid, not sorted magnitudes}

The arbitrary-node argument forgets the canonical coordinates and retains only the sorted
integers $m_i$.  Equation \eqref{eq:two-dimensional-grid} shows that rectangular subsets of
rows and columns produce structured bivariate exponential matrices, often Kronecker-like.
This is exactly the balanced $2\times2$ setting where ordinary six-exponentials estimates stop.
Integrality, positivity, and the one-depth-zero normalization may create a margin absent from
the general four-exponentials problem.

The right experiment is not another large generic determinant.  It is an optimized determinant
or minor selected from triangular coordinate regions
\[
 n+k\beta\le T,
 \qquad a+b\thetaq\le L,
\]
with exact accounting of archimedean size, zero estimates, and $2$-/$3$-adic valuations.
Computer algebra can search small shapes for unexpectedly large systematic factors before a
general theorem is attempted.

\subsection{Priority 3: denominator-sensitive counting on the compact arc}

The compact orbit formulation asks for infinitely many points on
$y=x^{\thetaq}$ in $[1,2)\times[1,3)$ whose first denominator is a power of $2$, whose second
denominator is a power of $3$, and whose denominator exponents are synchronized by the same
integer $\lfloor k\beta\rfloor$.

Standard rational-point counting treats denominator size but usually not this support and
synchronization.  A theorem of the following flavor would suffice:
\begin{quote}
On a compact transcendental arc with nonzero curvature, the number of points of logarithmic
height at most $H$ whose coordinate denominators are supported on two prescribed disjoint
prime sets and have a common exponent parameter is $o(H)$.
\end{quote}
The failed-conjecture orbit has $\asymp H$ such points.  Butler's connection between Wilkie-type
statements and real exponential conjectures suggests that this direction is genuinely close
to the transcendence frontier rather than a routine application of Pila--Wilkie.

\subsection{Priority 4: mixed archimedean--radical invariants}

The normalized radical $E=w^{\thetaq}$ has exact degree $d$ and binomial minimal polynomial.
Promising questions include:

\begin{itemize}
\item Can the logarithmic relation
$\log E\log2-\log w\log3=0$ force an impossible regulator or height relation in the Kummer
field $\Q(E,\zeta_d)$?
\item Does simultaneous normalization at both outputs constrain $d$ and $e$ through field
intersections or discriminants beyond the known gcd condition?
\item Can one use all complex branches
$\log(\zeta_d^jE)=\log E+2\pi i j/d$ in a multi-logarithm determinant whose algebraic entries
have controlled height?
\item Is there a useful $p$-adic logarithm branch compatible with the real identity after
raising to the normalized degree?
\end{itemize}

A warning is essential: merely taking norms or conjugates of $X^d-c/3^a$ reproduces known
relations.  The new invariant must mix conjugation with the special logarithmic direction.

\subsection{Priority 5: force a third algebraic exponential}

The denominator-controlled theorem in \path{RationalThirdBase.lean} gives a precise gateway.
For an auxiliary integer $q$ containing a prime outside $\{2,3\}$, it is enough in several
forms to prove that $q^x$ is rational with denominator supported by the known outputs.  Then
monomial injectivity and the three-base determinant force $x\in\Q$, hence $x\in\Z$.

Research should therefore search for natural expressions in $M,A$ whose logarithms introduce
a new prime direction but whose $x$th powers simplify.  Straight monomials cannot work;
possible candidates would involve norms, resultants, or algebraic combinations arising from
the primitive radical field.

\subsection{Priority 6: computational reconnaissance with exact certificates}

The finite scan through $2^{20}$ is evidence, not a path to infinity.  Computation can still be
valuable in three focused roles:

\begin{enumerate}
\item search candidate determinant shapes for large systematic $2$- or $3$-adic valuations;
\item enumerate small primitive radical patterns $(w,d,a,c)$ and record local obstructions,
with exact rational inequalities rather than floating-point rejection;
\item compress finite-check certificates through continued-fraction tiers or reusable interval
lemmas so that larger ranges can be replayed by the Lean kernel.
\end{enumerate}

The goal should be pattern discovery and theorem generation, not merely extending a decimal
barrier.

\section{Lean formalization blueprint for the new results}
\label{sec:lean-blueprint}

The new arguments are deliberately organized into reusable lemmas.  A possible module layout
under
\lean{ExponentialIdentities/TwoBaseIntegerExponent/}
is as follows.

\subsection{\lean{ArbitraryNodeDeterminant.lean}}

This module should formalize the ordinary-power determinant and Theorem~\ref{thm:arbitrary-node}.
The main dependencies are:

\begin{enumerate}
\item a determinant/divided-difference identity for the matrix with columns
$1,x,\ldots,x^{r-1},f(x)$;
\item a generalized mean-value theorem for divided differences at distinct ordered nodes;
\item the existing derivative tower
\lean{rpowDerivativeFamily} and
\lean{fallingRpowCoeff} from \path{HigherDifference.lean};
\item the integer-entry determinant lower bound already used elsewhere in the project.
\end{enumerate}

A target statement, schematically, is:

\begin{lstlisting}[style=leanstyle,caption={Proposed arbitrary-node theorem signature.}]
theorem vandermonde_lower_bound_of_twoBaseNaturalCandidates
    {r : N} (hr : 2 <= r)
    (m : Fin (r + 1) -> N)
    (hstrict : StrictMono m)
    (hcand : forall i, TwoBaseNaturalCandidate (m i)) :
    (r.factorial : R) / |fallingRpowCoeff logThreeDivLogTwo r| *
        (m 0 : R) ^ ((r : R) - logThreeDivLogTwo)
      <= productOverPairs fun i j => ((m j : R) - m i)
\end{lstlisting}

It may be easier first to state the theorem for a strictly increasing \lean{Finset} or a
\lean{Fin}-indexed sequence with an explicit pair product.

\subsection{\lean{GeneralizedVandermonde.lean}}

The positivity theorem is useful beyond this conjecture.  Two formalization routes are
available.

\begin{description}
\item[Wronskian route.]
Prove the determinant formula \eqref{eq:wronskian}; invoke or develop the criterion that
positive initial Wronskians imply an extended complete Chebyshev system and positive ordered
evaluation determinants.

\item[Zero-count route.]
After the substitution $x=e^t$, prove by induction that a nonzero exponential polynomial with
$r+1$ distinct real exponents has at most $r$ zeros counting multiplicity.  Deduce nonvanishing
of the evaluation determinant, then determine its sign by continuity and a coalescing-node
limit.
\end{description}

The zero-count route may require more real-analysis infrastructure, while the Wronskian route
may require adding a small Chebyshev-system API.  For an initial project-specific proof, one
can prove only nonvanishing and then compare signs through total positivity of the exponential
kernel $e^{st}$.

\subsection{\lean{MixedExponentSemigroup.lean}}

Define a pair-indexed exponent
\begin{lstlisting}[style=leanstyle]
def mixedExponent (ab : N x N) : R :=
  ab.1 + ab.2 * logThreeDivLogTwo
\end{lstlisting}
and prove:

\begin{itemize}
\item injectivity from irrationality of $\thetaq$;
\item integrality of
$m^{a+b\thetaq}=m^a(m^{\thetaq})^b$ for candidates;
\item positivity and integer-valuedness of the corresponding generalized determinant;
\item finite ordering of any chosen list of pairs by their real exponents.
\end{itemize}

For the quantitative theorem, it is unnecessary initially to define the globally sorted
semigroup.  One may accept an arbitrary strictly increasing finite exponent list and a proof
that every exponent has a pair representation.

\subsection{\lean{SemigroupDividedDifferenceBound.lean}}

Formalize Lemma~\ref{lem:falling-bound} and the strip bound
\eqref{eq:power-strip-bound}.  The main analytic statement should be a bound on every Newton
row:
\begin{lstlisting}[style=leanstyle]
abs (dividedDiff (fun x : R => x ^ lambda) nodesPrefix)
  <= 2 ^ (2 * lambda + i) * N ^ (lambda - i)
\end{lstlisting}
under $N\le x\le2N$ and $\lambda\ge0$.

The determinant estimate then follows from \lean{Matrix.det_apply} or a suitable
Leibniz-bound lemma.  It would be worth adding a generic API:
\begin{quote}
if every entry in row $i$, column $j$ is bounded by a separable quantity
$A_iB_j$, then the determinant is at most $(r+1)!\prod_iA_i\prod_jB_j$.
\end{quote}

\subsection{\lean{SemigroupCounting.lean}}

For the explicit logarithmic-square corollary, the full asymptotic
\eqref{eq:semigroup-counting} is optional.  Formalize the square-box construction from
Lemma~\ref{lem:crude-semigroup}:

\begin{itemize}
\item the $(q+1)^2$ pairs in $[0,q]^2$ give distinct mixed exponents;
\item every such exponent is below $3q$ using $1<\thetaq<2$;
\item hence the $r$th order statistic is at most $6\sqrt r$;
\item sum and divide to obtain $\Sigma_r/R_r\le12/\sqrt r$.
\end{itemize}

Avoiding real-valued order statistics may be easier if the theorem chooses any $r+1$ pairs
from a square box and bounds their total weight directly.  For example, take the first
$r+1$ pairs from a fixed enumeration of $[0,q]^2$; after sorting by exponent, their sum is at
most $3q(r+1)$.  This proves the same coarse determinant exponent without constructing the
canonical ordered semigroup.

\subsection{\lean{SemigroupDyadicBound.lean}}

Combine the span theorem with the consecutive-window summation argument.  The final theorem
could state:

\begin{lstlisting}[style=leanstyle,caption={Proposed explicit dyadic bound.}]
theorem card_twoBaseNaturalCandidates_Icc_le_log_sq (X : N) (hX : 1 <= X) :
  ((Finset.Icc X (2 * X)).filter TwoBaseNaturalCandidate).card
    <= Nat.ceil (10000 * max 1 (Real.log X) ^ 2)
\end{lstlisting}

A cleaner all-natural statement may replace the real logarithm by the binary ceiling logarithm,
obtaining a larger but easier constant times \lean{Nat.clog 2 X ^ 2}.  Since the repository
already has logarithmic-square counting infrastructure in \path{MultiplicativeRank.lean}, its
finite-set definitions and cardinality lemmas can be reused while keeping the proof mechanism
independent.

\subsection{Recommended proof-audit discipline}

Each new headline theorem should be included in the aggregate import and checked with
\texttt{\#print axioms}.  The report status should remain \statusnew{} until the source is
committed and a full build succeeds.  Numerical constants in tables are explanatory only;
the formal theorem should use exact expressions or safe rational constants such as $1/24$ and
$10000$.

\section{A more detailed comparison with existing ProveIt methods}
\label{sec:comparison}

\subsection{Prime-valuation rank versus interpolation geometry}

The module \path{MultiplicativeRank.lean} proves that any three candidates have linearly
dependent prime-exponent vectors of rank at most two.  It does so by a three-base/six-
exponentials determinant: if three candidate bases were multiplicatively independent, the
integrality of all three $\thetaq$th powers would force $\thetaq$ rational, contradicting its
irrationality.  A finite candidate set therefore lives in a rational plane in the free vector
space on primes.  Choosing two prime coordinates injectively encodes the set and yields an
explicit logarithmic-square count.

The semigroup determinant proves a logarithmic-square dyadic bound by a different mechanism:

\begin{center}
\begin{tabular}{p{0.42\textwidth}p{0.42\textwidth}}
\toprule
Prime-valuation method & Semigroup determinant method \\
\midrule
Uses prime factorization of candidate bases & Uses only integer values on $y=x^{\thetaq}$ \\
Global multiplicative rank at most two & Local total positivity and interpolation \\
Invokes a six-exponentials consequence & Uses elementary determinant integrality once
$\thetaq$ is fixed \\
Produces coordinates in two prime directions & Produces repulsion for arbitrary ordered nodes \\
Naturally global in magnitude & Naturally local in additive intervals \\
Already kernel-verified & Paper proof proposed for formalization \\
\bottomrule
\end{tabular}
\end{center}

Their agreement at logarithmic-square scale is suggestive.  Both see a two-dimensional
structure: prime-exponent rank two on one side, mixed-exponent lattice dimension two on the
other.  A future proof may need to combine these dual lattices rather than refine either in
isolation.

\subsection{Exact conditional rank one over a two-generator monoid}

Conditional on failure, the candidate monoid is not merely rank at most two; it is exactly
free on $2$ and $M$.  In prime-valuation space its points form the nonnegative lattice generated
by the vectors of $2$ and $M$.  In mixed-exponent evaluation space, the values factor as in
\eqref{eq:two-dimensional-grid}.  This exact grid is much stronger than the unconditional
rank theorem, but all presently known consequences remain compatible with it.

The strongest likely synthesis is a determinant whose rows use primitive-generator
coordinates and whose columns use mixed exponents.  Such a matrix is simultaneously:

\begin{itemize}
\item integral;
\item totally positive after ordering by magnitude;
\item explicitly factored into powers of $2,3,M,A$;
\item constrained by affine primitivity and simultaneous output normalization.
\end{itemize}

No current module appears to combine all four properties.

\subsection{Finite differences as confluent determinants}

Forward differences correspond to the equally spaced specialization of interpolation, while
Wronskians correspond to confluent nodes.  The new determinant hierarchy therefore provides a
single conceptual envelope for:

\begin{itemize}
\item second-, third-, and arbitrary-order finite differences;
\item arbitrary-node ordinary-power determinants;
\item mixed-monomial generalized Vandermonde determinants;
\item possible confluent determinants using derivatives or repeated coordinate directions.
\end{itemize}

A useful formalization strategy is to build a general interpolation-determinant layer, then
recover \path{HigherDifference.lean} as a specialization.  This would reduce duplicated
analysis and make experimental determinant choices easier to certify.

\section{Conclusions}
\label{sec:conclusion}

The Alaoglu--Erd\H{o}s conjecture remains unresolved here.  That is not for lack of structural
information.  Conditional on one counterexample, the pinned ProveIt corpus determines the
entire solution set, candidate monoid, exact block counts, asymptotic population, rotation
statistics, primitive odd core, rational-power index, radical degree, and simultaneous output
normalization.  The surviving obstruction is exceptionally concentrated:
\[
 \text{exclude }w^{\log_2 3}\in\Z[1/3]
 \text{ for every primitive odd }w>1.
\]

The new arbitrary-node determinant proves that candidate lattice points cannot cluster freely;
all Vandermonde products obey a power-law lower bound.  The mixed-exponent semigroup then gives
arbitrarily high-order positive integer determinants and a repulsion exponent tending to one.
It yields a second, conceptually independent logarithmic-square counting theorem.

The most valuable outcome is the quantified failure point.  Pure archimedean interpolation
requires about $(\log X)^2$ nodes to operate at scale $X$, whereas the exact conditional monoid
provides only $\log X$.  This makes future work more focused.  Improving constants or extending
finite searches will not close the gap.  The plausible routes are those capable of supplying
one additional logarithm of determinant strength: $2$- or $3$-adic divisibility, a
synchronized-denominator counting theorem on the compact exponential arc, or a mechanism that
forces a third algebraically independent exponential direction.

The new results are sufficiently elementary and modular to formalize now.  Even if the full
conjecture ultimately requires a major transcendence theorem, adding arbitrary-node and
semigroup-determinant APIs to ProveIt would sharpen the machine-checked frontier and provide the
right experimental framework for the next attack.

\appendix

\section{Detailed proof of the generalized Vandermonde sign}
\label{app:gv-proof}

This appendix expands Lemma~\ref{lem:generalized-vandermonde} to make its analytic input fully
transparent.

Let $0\le\lambda_0<\cdots<\lambda_r$ and define
$f_j(x)=x^{\lambda_j}$ on $(0,\infty)$.  For $k\le r$,
\[
 f_j^{(k)}(x)=\fall{\lambda_j}{k}x^{\lambda_j-k}.
\]
The $k$th initial Wronskian is
\begin{align*}
 W_k(x)
 &=\det\bigl(f_j^{(i)}(x)\bigr)_{0\le i,j\le k}\\
 &=x^{\sum_{j=0}^{k}\lambda_j-k(k+1)/2}
   \det\bigl(\fall{\lambda_j}{i}\bigr)_{0\le i,j\le k}.
\end{align*}
The polynomial $\fall{z}{i}$ is monic of degree $i$.  Replacing the monomial basis
$1,z,\ldots,z^k$ by the falling-factorial basis is unit lower-triangular.  Therefore
\[
 \det\bigl(\fall{\lambda_j}{i}\bigr)_{i,j=0}^{k}
 =\prod_{0\le i<j\le k}(\lambda_j-\lambda_i)>0.
\]
Hence every initial Wronskian is positive.

One standard characterization of extended complete Chebyshev systems now implies that for
$x_0<\cdots<x_r$,
\[
 \det(f_j(x_i))_{i,j=0}^{r}>0.
\]
To avoid using that characterization as a black box, proceed inductively.  Suppose a nontrivial
combination
\[
 F(x)=\sum_{j=0}^{r}c_jx^{\lambda_j}
\]
has $r+1$ distinct positive zeros.  Put $x=e^t$ and factor the lowest exponential:
\[
 e^{-\lambda_0t}F(e^t)
 =c_0+\sum_{j=1}^{r}c_je^{(\lambda_j-\lambda_0)t}.
\]
Rolle's theorem gives at least $r$ zeros of its derivative, an exponential polynomial with at
most $r$ terms.  Induction gives a contradiction unless all coefficients vanish.  Thus the
evaluation determinant is nonzero.  Its sign cannot change on the connected chamber
$x_0<\cdots<x_r$.  Let the nodes coalesce as $x_i=x+i\varepsilon$ and divide by the positive
ordinary Vandermonde.  As $\varepsilon\downarrow0$, the quotient tends to
$W_r(x)/\prod_{k=0}^{r}k!>0$.  Hence the original determinant is positive everywhere in the
ordered chamber.

\section{Numerical constants for arbitrary-node repulsion}
\label{app:constants}

For reference,
\[
 \thetaq=1.58496250072115618145\ldots,
 \qquad
 \thetaq(\thetaq-1)=0.92714362797110482756\ldots.
\]
The three-point span constant is
\[
 \left(\frac{8}{\thetaq(\thetaq-1)}\right)^{1/3}
 =2.051072395661874\ldots.
\]
For the ordinary-power order $r$, the raw span constant and exponent are
\[
 C_r=\left(\frac{r!}{|\fall{\thetaq}{r}|}\right)^{2/(r(r+1))},
 \qquad
 \alpha_r=\frac{2(r-\thetaq)}{r(r+1)}.
\]
The continuous critical point for $\alpha_r$ solves
\[
 -r^2+2\thetaq r+\thetaq=0,
\]
so
\[
 r_*=\thetaq+\sqrt{\thetaq^2+\thetaq}=3.6090841940\ldots,
\]
confirming that $r=4$ is the optimal integer order.

For the semigroup hierarchy, the first exponents and their pair labels are:
\begin{center}
\small
\begin{tabular}{rcl}
\toprule
Index & Pair $(a,b)$ & Exponent $a+b\thetaq$ \\
\midrule
0 & $(0,0)$ & $0$ \\
1 & $(1,0)$ & $1$ \\
2 & $(0,1)$ & $\thetaq$ \\
3 & $(2,0)$ & $2$ \\
4 & $(1,1)$ & $1+\thetaq$ \\
5 & $(3,0)$ & $3$ \\
6 & $(0,2)$ & $2\thetaq$ \\
7 & $(2,1)$ & $2+\thetaq$ \\
8 & $(4,0)$ & $4$ \\
9 & $(1,2)$ & $1+2\thetaq$ \\
10 & $(3,1)$ & $3+\thetaq$ \\
11 & $(0,3)$ & $3\thetaq$ \\
12 & $(5,0)$ & $5$ \\
\bottomrule
\end{tabular}
\end{center}

\section{Repository theorem inventory used in this report}
\label{app:inventory}

The following modules on the aggregate import surface were especially relevant.

\begin{longtable}{>{\raggedright\arraybackslash}p{0.34\textwidth}
                  >{\raggedright\arraybackslash}p{0.58\textwidth}}
\toprule
Module & Role \\
\midrule
\endhead
\path{TwoBaseIntegerExponent.lean}
& Basic definitions, fractional-part gaps, candidate exponent, and low-order curvature. \\
\path{FiniteCheck4096.lean}
& Exact natural-arithmetic certificates excluding nonintegral candidates below $4096$. \\
\path{Transcendence.lean}
& Gelfond--Schneider dichotomy; nonintegral solutions and $\thetaq$ are transcendental. \\
\path{FourExponentialsReduction.lean}
& Rational linear independence of $(1,x)$ and $(\log2,\log3)$; four exponentials implies the conjecture. \\
\path{Localization.lean}
& Equivalence with $u^{\thetaq}\notin\Z[1/3]$ for odd $u>1$. \\
\path{PrimitiveRadical.lean}
& Reduction to power-primitive odd bases. \\
\path{RadicalDegree.lean}
& Exact irreducible binomial and algebraic degree for the least rational power. \\
\path{PrimitiveOddCore.lean}
& Canonical common odd core of all hypothetical candidates. \\
\path{RationalPowerIndex.lean}
& Cyclic subgroup of rational powers and least index. \\
\path{PrimitiveGenerator.lean}
& Canonical least generator and exact free solution/candidate monoids. \\
\path{ExactCounting.lean}
& Exact unit-block and dyadic-block parametrization and cardinalities. \\
\path{QuadraticAsymptotic.lean}
& Sharp cumulative floor-sum bounds, quadratic limit, and growth equivalence. \\
\path{BlockWeyl.lean}
& Vanishing Fourier modes for exact block orbit segments. \\
\path{OutputNormalization.lean}
& Simultaneous primitive output decompositions and gcd constraints. \\
\path{RationalThirdBase.lean}
& Third-base rationality with controlled denominator support. \\
\path{ThreeSmooth.lean}
& Classification of auxiliary bases supported on $2$ and $3$. \\
\path{HigherDifference.lean}
& Arbitrary-order forward-difference mean value and progression obstruction. \\
\path{MultiplicativeRank.lean}
& Rank at most two for candidate prime-exponent vectors and logarithmic-square counting. \\
\bottomrule
\end{longtable}

\section{Suggested experimental determinant families}
\label{app:experiments}

Before developing a full $p$-adic theory, several finite experiments can test whether the
required amplification exists.

\subsection{Rectangular row and column boxes}

Choose rows $(n,k)$ with $0\le n<N_0$, $0\le k<K_0$ and columns $(a,b)$ with
$0\le a<A_0$, $0\le b<B_0$.  The entry is
\[
 (2^a3^b)^n(M^aA^b)^k.
\]
For complete rectangles, the matrix is a rearranged Kronecker product only when the row and
column cardinalities match in a compatible way.  Its determinant or maximal minors can be
factored symbolically for small boxes.  Record valuations at $2$ and $3$ as polynomials in
$a_0=\vpp{2}{M}$ and $b_0=\vpp{3}{A}$.

\subsection{Triangular boxes adapted to height}

Use
\[
 n+k\beta\le T,
 \qquad
 a+b\thetaq\le L.
\]
These are the natural row and column shapes from the exact counting theorem and mixed-exponent
semigroup.  Search for square minors with minimal archimedean weight and maximal forced
valuation.  The linear boundary slopes are irrational, so floor-sum fluctuations may prevent
cancellation patterns present in rectangles.

\subsection{Finite-difference row operations}

Order rows by $n$ for fixed $k$, or by $k$ for fixed $n$, and apply repeated difference
operators.  Differences introduce factors such as
\[
 2^a3^b-1,
 \qquad
 M^aA^b-1.
\]
These may have systematic $2$- or $3$-adic valuations for carefully selected column pairs.
Lifting-the-exponent lemmas can then provide exact lower bounds.  The risk is that the
archimedean size of the same factors grows too quickly; experiments should track both sides.

\subsection{Confluent columns}

Instead of only distinct mixed monomials, consider derivatives with respect to an auxiliary
continuous exponent parameter before specializing.  Confluent generalized Vandermonde
matrices introduce logarithmic factors, which are not integral, but common-denominator or
linear-form constructions may restore arithmetic control.  This is closer to Laurent's
interpolation determinant method and may reduce total exponent weight.

\section{Bibliographic and repository references}
\label{app:references}

\begin{thebibliography}{99}

\bibitem{AE1944}
L.~Alaoglu and P.~Erd\H{o}s,
\newblock \emph{On highly composite and similar numbers},
\newblock Transactions of the American Mathematical Society \textbf{56} (1944), 448--469.
\newblock DOI: \href{https://doi.org/10.1090/S0002-9947-1944-0011087-2}{10.1090/S0002-9947-1944-0011087-2}.

\bibitem{Baker1975}
A.~Baker,
\newblock \emph{Transcendental Number Theory},
\newblock Cambridge University Press, 1975.

\bibitem{Butler2017}
L.~A.~Butler,
\newblock \emph{A Diophantine approach to the three and four exponentials conjectures},
\newblock Ramanujan Journal \textbf{42}(1) (2017), 199--221.
\newblock DOI: \href{https://doi.org/10.1007/s11139-015-9728-2}{10.1007/s11139-015-9728-2}.

\bibitem{KarlinStudden}
S.~Karlin and W.~J.~Studden,
\newblock \emph{Tchebycheff Systems: With Applications in Analysis and Statistics},
\newblock Interscience, 1966.

\bibitem{KuipersNiederreiter}
L.~Kuipers and H.~Niederreiter,
\newblock \emph{Uniform Distribution of Sequences},
\newblock Wiley, 1974; Dover reprint, 2006.

\bibitem{Roy1992}
D.~Roy,
\newblock \emph{Matrices whose coefficients are linear forms in logarithms},
\newblock Journal of Number Theory \textbf{41}(1) (1992), 22--47.
\newblock DOI: \href{https://doi.org/10.1016/0022-314X(92)90081-Y}{10.1016/0022-314X(92)90081-Y}.

\bibitem{Schneider1957}
Th.~Schneider,
\newblock \emph{Einf\"uhrung in die transzendenten Zahlen},
\newblock Springer, 1957.

\bibitem{Waldschmidt2000}
M.~Waldschmidt,
\newblock \emph{Diophantine Approximation on Linear Algebraic Groups: Transcendence
Properties of the Exponential Function in Several Variables},
\newblock Grundlehren der mathematischen Wissenschaften 326, Springer, 2000.
\newblock DOI: \href{https://doi.org/10.1007/978-3-662-11569-5}{10.1007/978-3-662-11569-5}.

\bibitem{Waldschmidt2005}
M.~Waldschmidt,
\newblock \emph{Variations on the Six Exponentials Theorem},
\newblock in \emph{Algebra and Number Theory}, pp.~338--355, 2005.
\newblock DOI: \href{https://doi.org/10.1007/978-93-86279-23-1_23}{10.1007/978-93-86279-23-1\_23}.

\bibitem{ProveIt}
V.~Reshetnikov,
\newblock \emph{ProveIt repository},
\newblock \url{https://github.com/VladimirReshetnikov/ProveIt},
\newblock pinned in this report at commit
\commit{5d3982ceb12644fdd3499569d9db5f62f9a31dc6}.

\bibitem{CombinedReport}
The ProveIt project,
\newblock \emph{The Alaoglu--Erd\H{o}s Integer-Exponent Conjecture: Machine-Checked Partial
Results, Exact Conditional Structure, and Research Directions},
\newblock source file
\path{Analysis/ExponentialIdentities/Article/alaoglu-erdos-combined-report/alaoglu-erdos-combined-report.tex},
August 20, 2026.

\end{thebibliography}

\end{document}
